\section{The R\'enyi bounds: full proofs and further results}\label{sec:sup-renyi}
The intermediate-order R\'enyi proof works with a \emph{relaxation} of the
quantum problem. Take three arbitrary four-point probability vectors
$\lambda^{(j)}$, $j=2,3,4$, and set
$P=\sum_{j,i}(\lambda_i^{(j)})^2$ and
$T=\sum_{j,i}(\lambda_i^{(j)})^3$. Impose the necessary constraints
$P\ge1$, $F=17P-12T-12\ge0$, and the quadratic purity-variance bound
(QK). We do not assume these inequalities characterise physically
realisable cut spectra. An upper bound for this larger set proves the
quantum bound; identifying the equality states requires the separate
rigidity arguments of the paper. The bound (QK)
is deduced below from the sharper six-copy certificate (QV).

The remaining optimisation has two anchor orders. Determinant monotonicity
reduces individual spectra to repeated-root charts; the supporting bounds
then cover an HS region, a strict intermediate band, a capped region and a
narrow cluster region. Single crossing transfers the HS anchor to lower
orders, and monotonicity of R\'enyi entropy transfers the cluster anchor to
higher orders. The final subsections give independent estimates and explain
why local filtering cannot replace the unrestricted compatibility argument.

\subsection*{Proofs of results of the paper}
\subsection{Proof of Theorem~\ref{thm:fqe-sharp-quadratic-variance} and Corollary~\ref{thm:fqe-quadratic-variance}}\label{sec:sup-qk}
\begin{restated}[Sharp quadratic purity variance; exact-rational certificate, QVA]{Theorem}{thm:fqe-sharp-quadratic-variance}
Every unit pure four-qubit state satisfies
\[
\boxed{\qquad F^2+3P(P-1)\ge6V_P,
\qquad V_P=\sum_{j=2}^4(p_j-P/3)^2.\qquad}
\tag{QV}
\]
Equality holds along the whole cluster-to-Bell family \eqref{eq:fqe-quadratic-family}, so the
coefficient $3$ cannot be decreased.
\end{restated}
Write $p_j=\operatorname{tr}\rho_{1j}^2\in[1/4,1]$ and $U=\sum_jp_j^2$. Maximising the convex function $U$ at fixed
$P$ over the box $[1/4,1]^3$, filling one coordinate at a time towards an endpoint, gives the \emph{purity box}
\[
U\le
\begin{cases}
(P-1/2)^2+1/8,&1\le P\le3/2,\\
1+(P-5/4)^2+1/16,&3/2\le P\le9/4,\\
2+(P-2)^2,&9/4\le P\le3.
\end{cases}
\tag{JB}\label{eq:fqe-purity-box}
\]
The family \eqref{eq:fqe-quadratic-family} has $P=1+r^2/2$, $F=(1-r^2)/2$ and $V_P=(1+2r^2)^2/24$, which gives
equality and, for $r>0$, $(6V_P-F^2)/(P(P-1))=3$ exactly (QV-3). Proposition~\ref{prop:fqe-critical-quadratic-variance}
covers critical states, including $P=1$ by Lemma~\ref{lem:fqe-purity}. When $F\ge3/2-P$, the purity box
\eqref{eq:fqe-purity-box} gives the inequality directly by substituting into $F^2+5P^2-3P-6U$ (QV-2).
Thus only noncritical states with $1<P<3/2$ and $F<3/2-P$ need the certificate, which
treats all states at once.

\begin{proof}
For $z\in\mathbb C^{16}$ put $N=\lVert z\rVert^2$, use unnormalised marginals in $P,T,U$, and set
\[
 \mathfrak F=17NP-12T-12N^3,\qquad
 \mathfrak D=\mathfrak F^2+5N^2P^2-3N^4P-6N^2U.
\]
Then $\mathfrak D(z)=N^6\bigl(F^2+3P(P-1)-6V_P\bigr)(z/\sqrt N)$, so it suffices to prove $\mathfrak D\ge0$.
It is a local-unitary invariant of bidegree $(6,6)$, symmetric under qubit permutations and complex
conjugation, and it vanishes at HS and at the cluster state, which are equality cases. We exhibit it as a sum of
squared norms.

Put $\widetilde z=\epsilon^{\otimes4}\bar z$, where $\epsilon=\bigl(\begin{smallmatrix}0&1\\-1&0\end{smallmatrix}\bigr)$.
Since $\epsilon\bar U=U\epsilon$ for $U\in SU(2)$, apply the six-copy spin intertwiners to
$z^{\otimes a}\otimes\widetilde z^{\otimes(6-a)}$, $a\in\{3,4\}$, and stack all $24$ qubit permutations.
For a spin tuple $p=(j_1,\dots,j_4)$, one twice-spin $j_i\in\{0,2,4,6\}$ of the six copies at each qubit, call the resulting map $W_{a,p}$ and equip its magnetic rows with the positive weights
\[
 w_p=\bigotimes_{j\in p}\left(\frac{\operatorname{lcm}_{r}\binom jr}{\binom jk}\right)_{k=0}^{j},
\]
repeated over the permutations. These are the invariant spin inner products: the spin-$j/2$ lowering basis has weights $1/\binom jk$.
The data in the verifier map specifies integer multiplicity vectors $v_{ki}$, integer weights $q_{ki}$
and an integer denominator $D$. Put $\Phi_{ki}=W_{a_k,p_k}v_{ki}$.
Intertwining makes their squared norms local-unitary invariant; summing over permutations makes them
qubit-symmetric, and the real intertwiners and weights make them conjugation-even.

There are three finite checks: the highest-weight identities and vanishing of each $\Phi_{ki}$ at HS and
cluster (QVA-1); positivity of $D$ and every $q_{ki}$ (QVA-2); and the identity
\begin{equation}\label{eq:fqe-qv-allstate-sos}
 2D\,\mathfrak D(z)=\sum_{k,i}q_{ki}\lVert\Phi_{ki}(z)\rVert_{w_{p_k}}^2
\end{equation}
at the $83$ determining states (QVA-3).
Both sides lie in the space $\mathcal V$ of qubit-symmetric, conjugation-even local-unitary
invariants of bidegree $(6,6)$ vanishing at HS and cluster. The ambient space has dimension $85$
(Section~\ref{sec:fqe-invariant-dimensions}, SD-4); HS and cluster evaluations are independent,
since the normalised values of $(N^6,N^3T)$ are $(1,5/12)$ and $(1,3/8)$. Thus $\dim\mathcal V=83$.
QG supplies $83$ Gram entries in $\mathcal V$, built from five-copy covariants multiplied by $N$
or coupled to one further copy (QG-1), whose evaluation matrix at these states
has determinant $627608$ modulo $1000003$ (QG-3). Evaluation is therefore injective on $\mathcal V$,
and \eqref{eq:fqe-qv-allstate-sos} holds identically. Every summand is nonnegative, so $\mathfrak D\ge0$.
\end{proof}

\begin{restated}[Quadratic purity variance]{Corollary}{thm:fqe-quadratic-variance}
Every unit pure four-qubit state satisfies
\[
\boxed{\qquad F^2+\frac92(P-1)\ge6V_P.\qquad}
\tag{QK}
\]
Equality holds exactly for HS cut spectra, cluster cut spectra, or products
of two Bell pairs, up to ordering of cuts and eigenvalues.
The coefficient $9/2$ cannot be decreased, and the coefficient $6$ cannot
be increased while retaining $9/2$.
\end{restated}
\begin{proof}
Write $\mathcal K=F^2+\frac92(P-1)-6V_P$ and $\mathcal V_Q=F^2+3P(P-1)-6V_P$, so that
$\mathcal K-\mathcal V_Q=3(P-1)(3/2-P)$. By Lemma~\ref{lem:fqe-purity}, $P\ge1$. For $1\le P\le3/2$,
Theorem~\ref{thm:fqe-sharp-quadratic-variance} gives $\mathcal K\ge3(P-1)(3/2-P)\ge0$, strictly for $1<P<3/2$.
For $P\ge3/2$ the purity box \eqref{eq:fqe-purity-box} gives
$\mathcal K\ge F^2+4(P-\tfrac32)(\tfrac{27}8-P)$ on $[3/2,9/4]$ and
$\mathcal K\ge F^2+(P-\tfrac94)(\tfrac{39}2-4P)+\tfrac{27}8>0$ on $[9/4,3]$ (QV-2), so $\mathcal K\ge0$, strictly for
$P>3/2$. Equality therefore forces $P=1$ or $P=3/2$.

At $P=3/2$ the box gives $\mathcal K\ge F^2$, so equality forces $F=0$ and $U=9/8$; the purities are a permutation of
$(1,1/4,1/4)$. The marginal of purity one makes the state a product across that cut. On either crossed cut its
marginal is a product of two one-qubit states, and purity $1/4$ forces both to be maximally mixed, so both factors
are Bell pairs. At $P=1$, Proposition~\ref{prop:fqe-critical-variance} gives
$\mathcal K=F^2-6V_P\ge s^4(1-s)(2-s)$ with $0\le s\le1$ (QV-1). Equality requires $s\in\{0,1\}$ and equality in
\eqref{eq:fqe-critical-variance}, which gives exactly HS or cluster spectra. All three types attain equality.
Bell products force the coefficient $9/2$, and cluster spectra force the coefficient $6$.
\end{proof}

The Lean formalisation proves~\eqref{eq:fqe-qv-allstate-sos} for all $z$ directly, without the dimension count
(Section~\ref{sec:fqe-lean}). An older direct certificate for (QK), from eighty squared covariant norms in sixteen modules, is
kept as a regression (QG-2, QG-4).

\subsection{Proof of Lemma~\ref{lem:fqe-renyi-determinant}}\label{sec:sup-determinant}
At fixed trace, second moment and third moment, the determinant is the only
remaining elementary symmetric function of a four-point spectrum. The lemma
shows that increasing this determinant improves R\'enyi entropy for
$1<q<2$, reducing each spectral optimisation to a repeated-root boundary.
\begin{restated}[Determinant monotonicity at fixed second and third moments]{Lemma}{lem:fqe-renyi-determinant}
For $1<q<2$, let $\lambda,\lambda'$ be four-point spectra with the same second moment $p$ and third moment $t$.
If $\prod_i\lambda_i<\prod_i\lambda'_i$, then $H_q(\lambda)<H_q(\lambda')$.
Consequently, at fixed $p$ and $t$, the maximum of $H_q$ is attained at a spectrum of one of the two types
\[
(a,b,c,c)\quad\text{or}\quad(a,a,b,c),
\qquad a\ge b\ge c\ge0,
\tag{DM}
\]
including their degenerations.
\end{restated}
\begin{proof}
Newton's identities fix $e_1=1$, $e_2=(1-p)/2$ and $e_3=(1-3p+2t)/6$ (UR-3). Hence, for $s>0$ and $d=\prod_i\lambda_i$,
\[
 \prod_i(\lambda_i+s)=M(s)+d,\qquad \sum_i\prod_{j\ne i}(\lambda_j+s)=M'(s),\qquad M(s)=s^4+s^3+e_2s^2+e_3s,
\]
where $M$ depends only on $(p,t)$ and both quantities are positive. Since $x^2/(x+s)=x-s+s^2/(x+s)$,
\[
 \sum_i\frac{\lambda_i^2}{\lambda_i+s}=1-4s+\frac{s^2M'(s)}{M(s)+d},
\]
which is strictly decreasing in $d$ for every $s>0$. For $1<q<2$ and $x\ge0$, the substitution $s=xu$ and
$\int_0^\infty u^{q-2}(1+u)^{-1}\,du=\pi/\sin(\pi(q-1))$ give
\[
 x^q=\frac{\sin(\pi(q-1))}{\pi}\int_0^\infty s^{q-2}\,\frac{x^2}{x+s}\,ds .
\]
Each of these integrals converges, so summing over the four eigenvalues gives, for $d<d'$,
\[
 \sum_i\lambda_i^q-\sum_i\lambda_i'^q=\frac{\sin(\pi(q-1))}{\pi}\int_0^\infty s^{q}M'(s)
 \Bigl(\frac1{M(s)+d}-\frac1{M(s)+d'}\Bigr)ds>0,
\]
and $H_q=(1-q)^{-1}\log_2\sum_i\lambda_i^q$ gives $H_q(\lambda)<H_q(\lambda')$.

For (DM), the ordered spectra $\lambda_1\ge\lambda_2\ge\lambda_3\ge\lambda_4\ge0$ with trace one and the given $p,t$
form a compact set, so the determinant attains its maximum on it, and by the first part a spectrum of maximal
determinant maximises $H_q$. Suppose it has $\lambda_1>\lambda_2$ and $\lambda_3>\lambda_4$. Let
$g(z)=\prod_i(z-\lambda_i)$, $m_1=(\lambda_1+\lambda_2)/2$ and $m_2=(\lambda_3+\lambda_4)/2$. Then $g(m_1)<0$,
$g(m_2)<0$ and $g(\lambda_2)=g(\lambda_4)=0$, so for $0<\delta<\min\{-g(m_1),-g(m_2)\}$ the quartic $g+\delta$ changes
sign on each of $(\lambda_4,m_2)$, $(m_2,\lambda_2)$, $(\lambda_2,m_1)$ and $(m_1,\infty)$. Its four roots are therefore
distinct and exceed $\lambda_4\ge0$. They have the same $e_1,e_2,e_3$, hence the same $p,t$, and product $d+\delta$,
contradicting maximality. So $\lambda_1=\lambda_2$ or $\lambda_3=\lambda_4$, which is (DM).
\end{proof}

\subsection{The classical relaxation: multiplier table}\label{sec:sup-renyi-classical}
The proof of Theorem~\ref{thm:fqe-renyi-classical} of the paper uses the multiplier table below and the
supporting lemmas, at the orders $q_s=1587/1000$ and $q_u=8/5$ of that theorem.
The next subsection combines them; the supporting inequalities follow afterwards.

\paragraph{Multiplier table.}
The table specifies the supporting function $\Lambda_\theta$ of Lemma~\ref{lem:fqe-dual}, with
$\theta=(\alpha,\beta,\gamma)$ multiplying the centred cut purity $p-\frac13$, the moment slack $f-4$ and minus
the centred squared purity, $-(p^2-\frac19)$, respectively. Here $\varepsilon=P-1$, $F=17P-12T-12$, and $h_q$ is the mean R\'enyi entropy of $M_4$.
Only two rational multiplier rows are needed: L for the strict band and HT
for the cluster sliver, both at $q_s=1587/1000$. Divide the integer entries
by $10000$. The HS supporting inequality at $q_u=8/5$ has its own explicit
multiplier (Lemma~\ref{lem:fqe-hs-polynomial}); the cubic bounds use none.
The partition in Table~\ref{tab:fqe-renyi-coverage} specifies where each is used.
The L estimate is proved on the slightly larger rectangle
$[0,1/9]\times[3/20,11/25]$ in $(\varepsilon,F)$; the HT estimate uses
$\varepsilon\le1/100$ and distinguishes one cut of purity above $47/100$
from two of purity at most $29/100$.
\begin{table}[htbp]
\centering\small
\begin{tabular}{crrr}
\toprule
Band&$10000\alpha$&$10000\beta$&$10000\gamma$\\
\midrule
% CZ-ROWS-BEGIN
L&39700&1500&15900\\
% CZ-ROWS-END
\midrule
% HT-ROWS-BEGIN
HT&35425&1872&11235\\
% HT-ROWS-END
\bottomrule
\end{tabular}
\caption{Exact supporting multipliers at the lower anchor $q_s$; divide
by $10000$. L supplies a strict common bound, while HT uses separate bounds
for the distinguished cut and the two near-uniform cuts.}
\label{tab:fqe-dual-co}
\end{table}



\subsection{Proof of Theorem~\ref{thm:fqe-renyi-classical}}\label{sec:sup-renyi-classical-proof}
\begin{restated}[Two R\'enyi anchors]{Theorem}{thm:fqe-renyi-classical}
Let $q_s=1587/1000$ and $q_u=8/5$. Every triple satisfying $P\ge1$, $F\ge0$ and
\eqref{eq:fqe-quadratic-variance-proved} obeys one of the following.
\begin{enumerate}
\item If $F\ge3/20$ or $20\varepsilon+2F\ge1$, then $\mathcal E_{q_s}\le5/3$, with equality only at a permutation of
$(c_1,c_2,c_2)$.
\item Otherwise $\mathcal E_{q_u}\le h_{q_u}$, with equality only when all three spectra are HS spectra.
\end{enumerate}
\end{restated}
\begin{proof}
The following partition is exhaustive and disjoint. In the last four rows,
$20\varepsilon+2F<1$; in the last two, $F>11/25$.
\begin{table}[htbp]
\centering\small
\begin{tabular}{@{}p{.36\linewidth}cl@{}}
\toprule
Region (first applicable row) & Order & Bound used \\
\midrule
$20\varepsilon+2F\ge1$ & $q_s$ & Cubic envelope (BC) \\
$F<3/20$ & $q_u$ & HS support (HD) \\
$3/20\le F\le11/25$ & $q_s$ & Strict band L (LT, SHD, LHD) \\
$\max_jp_j\le47/100$ & $q_s$ & Capped cubic bound (CG) \\
$\max_jp_j>47/100$ & $q_s$ & Cluster sliver (NU, SL) \\
\bottomrule
\end{tabular}
\caption{The anchor partition. Rows are applied in order: each triple is handled by the first row whose region
contains it, so every row excludes the regions above it. HD gives HS equality; BC gives cluster equality; the other
three regions are strict. SHD and CG both use the monotonicity of $A_q$ in Lemma~\ref{lem:fqe-binary-quartic}, proved by Cauchy's formula;
SHD adds the quartic checks SHD-3.}
\label{tab:fqe-renyi-coverage}
\end{table}

\emph{Cubic envelope.} Corollary~\ref{cor:fqe-renyi-high-f}, with $q_s>q_b$,
handles the first row and classifies its equality triples.
Henceforth $\varepsilon<1/20$, $F<1/2$, and
\[
 p_j\le P-1/2<11/20<9/16<49/75.
\]
These are the purity caps required by HD and L; they are consequences of
the region, not hypotheses on physical states.

\emph{HS zone.} For $F<3/20$, Lemma~\ref{lem:fqe-hs-polynomial} gives
$\Lambda_\theta\le h_{q_u}$, strictly off HS, and
$\ell_\theta\ge0$ on $[0,1/9]\times[0,3/20]$, which contains $(\varepsilon,F)$.
The dual lemma therefore gives $\mathcal E_{q_u}\le h_{q_u}$,
with equality only at three HS spectra.

\emph{Strict band.} For $3/20\le F\le11/25$,
Lemma~\ref{lem:fqe-rational-dual} gives
$\Lambda_\theta<W-10^{-5}$, where $3W=5+\min\ell_\theta$ on the
L rectangle. Thus
\[
 3\mathcal E_{q_s}\le\sum_j\Lambda_\theta-\ell_\theta
 <5+\min\ell_\theta-\ell_\theta-3\cdot10^{-5}<5.
\]

\emph{Near the cluster triple.} The remaining region has $F>11/25$ and
$\varepsilon<(1-2F)/20<3/500<1/100$.
If every cut has $p_j\le47/100$, Corollary~\ref{cor:fqe-capped-cluster}
gives the strict ceiling. Otherwise label a cut with $p_2>47/100$.
It is unique, since two such cuts would give
$P>47/50+1/4>101/100$. The other two have
$p_j\le P-p_2-1/4<29/100$.
For this last case only, use determinant replacement
(Corollary~\ref{cor:fqe-renyi-six-variable}) to put each cut on an ordered
double-root chart, preserving every $p_j,t_j$ and the constraints.

\emph{Cluster sliver.} Set $X=(p_3-1/4)+(p_4-1/4)=1/2+\varepsilon-p_2\le1/25$.
Writing $x_j=p_j-1/4$, Lemma~\ref{lem:fqe-third-moment} for cuts $3,4$ gives
\[
 F\in I=[F_c-\kappa_aX,F_c+\kappa_aX],\qquad
 F_c=-1+8\varepsilon+9p_2-12t_2,\qquad \kappa_a=4\sqrt3/5,
\]
since the error is at most
$4\sqrt3(x_3^{3/2}+x_4^{3/2})\le4\sqrt3X^{3/2}\le\kappa_aX$.
The actual $F$ lies in $I\cap(11/25,1/2)$, exactly the window needed by SL.
Use the HT row $\theta_s=(3.5425,0.1872,1.1235)$ and put
\[
 \Lambda^\star_i=\Lambda_{\theta_s}(c_i),\qquad
 m=\min_{f\in\partial I}\ell_{\theta_s}(\varepsilon,f),\qquad
 \ell_0=\ell_{\theta_s}(0,1/2).
\]
NU bounds each of the two near-uniform cuts by $\Lambda^\star_2$.
SL bounds the distinguished cut by $\Lambda^\star_1+m-\ell_0$.
Concavity of $\ell_{\theta_s}$ in $F$ gives $\ell_{\theta_s}(\varepsilon,F)\ge m$; hence Lemma~\ref{lem:fqe-dual}
and the contact identity~\eqref{eq:fqe-dual-contact} give
\[
 3\mathcal E_{q_s}
 \le \Lambda^\star_1+2\Lambda^\star_2-\ell_0+m-\ell_{\theta_s}(\varepsilon,F)
 \le \Lambda^\star_1+2\Lambda^\star_2-\ell_0=5.
\]
Here NU and SL are Lemmas~\ref{lem:fqe-near-uniform-dual}
and~\ref{lem:fqe-cluster-sliver}.

\emph{Equality.} HD forces three HS spectra. The strict band and capped region are strict.
In the sliver, equality in SL forces the distinguished cut to be $c_1$ and $\varepsilon=0$.
Then $X=0$, so both other cuts are $c_2$, giving $F=1/2$, excluded in this region.
Thus the replaced triple, and hence the original triple, has a strict bound there.
Only the HS equality of HD and the cluster equality of the cubic envelope remain.
\end{proof}

\subsection{A monotone path from the repeated-largest chart}
\label{sec:fqe-low-tilt-reduction}
\begin{lemma}[Low third-moment tilt; analytic]\label{lem:fqe-low-tilt-reduction}
Let $3/2\le q\le8/5$ and $0\le\omega\le9/5$.
Among spectra $(a,a,b,c)$ with $a\ge b\ge c\ge0$ of a fixed purity,
$H_q-\omega t$ is maximised only by two equal pairs.
Consequently, whenever $0\le\beta\le3/20$, the maximum at fixed purity of the supporting function
$\Lambda_\theta$ of Lemma~\ref{lem:fqe-dual} is attained on the single pattern $(a,b,c,c)$.
\end{lemma}
\begin{proof}
\emph{The path.} Write a repeated-largest spectrum of purity $p$ as
\[
 (1/4+u,\,1/4+u,\,1/4-u+d,\,1/4-u-d),\qquad
 x=p-1/4=4u^2+2d^2,\quad 0\le d\le2u,
\]
the last condition being the ordering. Increase $u$ to $u_*=\sqrt x/2$ at fixed $x$, with
$d=\sqrt{(x-4u^2)/2}$. Then $d'=-2u/d$ and $d\le2u$, so the middle entry $1/4-u+d$ decreases,
stays at most $1/4+u$, and the smallest entry has derivative $2u/d-1\ge0$: it never falls below its
initial value $c\ge0$. The endpoint is two equal pairs. Along the path
\[
 t=\frac1{16}+\frac34x+12u^3-3xu,\qquad
 \frac{dt}{du}=36u^2-3x=6(4u^2-d^2)\ge0,
\]
strictly in the interior unless the path is constant.

\emph{The slope.} Put $r=q-1\in[1/2,3/5]$, $L=\ln2$ and $Z=\sum_i\lambda_i^q$, so that
$H_q=\log_2Z/(1-q)$. The path preserves trace and purity and moves the three distinct entries
$a,b,c$; interpolating $x^r$ quadratically at them gives $dZ=\frac q3[a,b,c]x^r\,dt$, hence
\[
 \frac{dH_q}{dt}=\eta(a,b,c):=-\frac{q[a,b,c]x^r}{3rLZ}.
\]
For distinct $a,b,c\ge0$ with mean $m$ and $0<r<1$,
\[
 -[a,b,c]x^r\ge\frac{r(1-r)}2\,m^{r-2}.
 \tag{DD}\label{eq:fqe-dd-power}
\]
Indeed, let $T$ be the cubic Taylor polynomial of $x^r$ at $m$. Then
$(x^r-T)''=-r(1-r)\{x^{r-2}-m^{r-2}-(r-2)m^{r-3}(x-m)\}\le0$ by the tangent inequality for the convex
function $x^{r-2}$, so $x^r-T$ is concave on $[0,\infty)$ and $[a,b,c](x^r-T)\le0$. As
$[a,b,c](x-m)^2=1$ and $[a,b,c](x-m)^3=a+b+c-3m=0$, $[a,b,c]T=-\frac12r(1-r)m^{r-2}$.
Here $m=(1-a)/3$. With $Z\le a^r\sum_i\lambda_i=a^r$ this gives
\[
 \eta(a,b,c)\ge\frac{q(1-r)3^{2-r}}{6L\,a^r(1-a)^{2-r}},\qquad
 a^r(1-a)^{2-r}=[a(1-a)]^r(1-a)^{2-2r}\le\frac9{16}\left(\frac49\right)^r\le\frac38,
\]
using $1/4\le a\le1/2$ and $r\ge1/2$. With $q(1-r)=q(2-q)\ge16/25$, $L<7/10$ and
$3^{2-r}\ge3^{7/5}>9/2$ (as $3^7>(9/2)^5$),
\[
 \eta(a,b,c)>\frac{64}{35}>\frac95\ge\omega.
\tag{LT}\label{eq:fqe-low-tilt-slope}
\]
Thus $d(H_q-\omega t)/du=(\eta-\omega)\,dt/du>0$ on every nonconstant path, and continuity covers zero or
coalescing entries. This proves the first claim.

For the second, Lemma~\ref{lem:fqe-renyi-determinant} reduces a fixed-purity maximum of $\Lambda_\theta$ to the
patterns $(a,b,c,c)$ and $(a,a,b,c)$. At fixed purity the variable part of $\Lambda_\theta$ is $H_q-12\beta t$ with
$12\beta\le9/5$, and two equal pairs already belong to the first pattern.
\end{proof}


\subsection{An analytic HS supporting inequality}\label{sec:fqe-hs-polynomial}

At the upper anchor $q_u=8/5$, choosing $\beta=1/12$ makes the
fixed-purity optimisation a bound on $H_{8/5}-t$. We first show that its
maximiser has one large eigenvalue, then prove the resulting one-variable
supporting inequality. The outer dual is nonnegative for $F\le3/20$,
which determines the split between the two anchor regions.
Throughout this subsection put $q=8/5$, $r=q-1=3/5$ and $L=\ln2$.

\begin{lemma}[A third-moment tilt at fixed purity; proved]\label{lem:fqe-hs-tilt}
Among four-point probability vectors of a fixed purity $p\le9/16$,
$H_{8/5}(\lambda)-\sum_i\lambda_i^3$ is maximised only by the one-large spectrum
\[
 \lambda(s)=\bigl((1+3s)/4,(1-s)/4,(1-s)/4,(1-s)/4\bigr),
 \qquad p=(1+3s^2)/4,
\]
up to permutation (uniform when $p=1/4$).
\end{lemma}
\begin{proof}
Let $Z=\sum_i\lambda_i^q$ and let $a,b,c$ be three distinct entries.
Their mean is at most $1/3$, so the divided-difference bound~\eqref{eq:fqe-dd-power} gives
$-[a,b,c]x^r\ge\frac12r(1-r)3^{2-r}$. Also
$Z=\sum_i\lambda_i\lambda_i^r\le p^r$ by concavity of $x^r$.
Consequently
\[
 \eta(a,b,c):=-\frac{q[a,b,c]x^r}{3rLZ}
 \ge\frac8{75L}\frac{3^{7/5}}{(9/16)^{3/5}}
 >\frac{16}{105}(3\cdot2^{12})^{1/5}>1.
 \tag{HTS}\label{eq:fqe-hs-tilt-slope}
\]
The last comparisons use $L<7/10$ and
$105^5=12762815625<3\cdot2^{32}=12884901888$.

Lemma~\ref{lem:fqe-low-tilt-reduction}, with $\omega=1$, reduces the maximum
to the pattern $(a,b,c,c)$. Write it as
\[
 (1/4+u+d,1/4+u-d,1/4-u,1/4-u),\qquad
 x=p-1/4=4u^2+2d^2,\quad 0\le d\le2u.
\]
Decrease $u$ to $u_0=\sqrt{x/12}$ at fixed $x$, setting
$d=\sqrt{(x-4u^2)/2}$. Ordering and positivity persist, and $b=c$ at
the endpoint: it is the one-large spectrum.
Along this path
\[
 t=\frac1{16}+\frac34x+3xu-12u^3,\qquad
 \frac{dt}{du}=3x-36u^2\le0.
\]
As in the preceding chart reduction, $dH_q/dt=\eta(a,b,c)>1$.
Thus decreasing $u$ strictly increases $H_q-t$ unless the spectrum
is already one-large. Continuity handles zero or coalescing entries.
At equality the strict paths force every determinant-maximal replacement to be
one-large. That replacement preserves $p,t$, so equality in the sharp third-moment bound
(Lemma~\ref{lem:fqe-third-moment}) forces the original spectrum to be the same.
The uniform case $p=1/4$ is immediate.
\end{proof}

\begin{lemma}[HS supporting inequality; proved]\label{lem:fqe-hs-polynomial}
At $q=q_u=8/5$ put $u=3^{3/5}$ and
\[
 k=\frac{16}{3L}\frac{u-1}{u+1},\qquad
 \beta=\frac1{12},\qquad \gamma=\frac{10}3,\qquad
 \alpha=\frac32k+\frac{65}{36}.
\]
Every four-point probability vector with $p\le9/16$ satisfies
$\Lambda_{(\alpha,\beta,\gamma)}(\lambda)\le h_{8/5}$,
with equality only at HS. For $0\le\varepsilon\le1/9$ and
$0\le F\le3/20$,
\[
 \ell_{(\alpha,\beta,\gamma)}(\varepsilon,F)
 \ge\frac{149}{324}\varepsilon
       +F\left(\frac1{12}-\frac59F\right)\ge0.
\]
\end{lemma}
\begin{proof}
At fixed $p$, the only variable part of $\Lambda$ is $H_q-t$.
Lemma~\ref{lem:fqe-hs-tilt} reduces its maximum to $\lambda(s)$.
Write $f(p)=H_q(\lambda(s))$, $p_* =1/3$, and
\[
 A(s)=\frac92(s-\tfrac13)^2(s+\tfrac16),\qquad
 W(p)=f(p)-\beta A(s)-\gamma(p-p_*)^2.
\]
Direct moment evaluation gives
\[
 \Lambda(\lambda(s))=W(p)+\frac32k(p-p_*),\qquad
 W(p_*)=h_q,\qquad W'(p_*)=-\frac32k,
 \qquad \frac{d^2 A}{dp^2}=\frac6s.
\]
We show that $W$ is first convex and then concave, allowing either
interval to be empty. The needed scalar calculation is explicit below.
Put
\[
 x=\left(\frac{1+3s}{1-s}\right)^{1/5},\quad
 J=\frac{L f''(p)}{2\gamma+6\beta/s},\quad
 h_4(x)=1+x+x^2+x^3+x^4.
\]
Since $p\le9/16$, $s<2/3$, so $1<x<2$ and $0<x^5-1<8$.
Differentiating
\[
 L\frac{dH_q(\lambda(s))}{ds}
 =-\frac{2(x^5+3)(x^3-1)}{x^8+3}
\]
gives the identity
\[
 \frac{dJ}{dx}
 =-\frac{8(x^5+3)^3 Q(x)}
 {5x^3(43x^5-31)^2(x^8+3)^3h_4(x)^3}.
 \tag{HCJ}\label{eq:fqe-hs-curvature-identity}
\]
Here, with $y=x-1$ and $R=x^5-1$,
$Q(x)=\sum_{j=0}^7R^jP_j(y)$, where the five columns below are the
coefficients of $1,y,y^2,y^3,y^4$ in $P_j$.
\[
\begin{array}{c|rrrrr}
 j&1&y&y^2&y^3&y^4\\\hline
% HD-CURVATURE-BEGIN
0&225600&457600&471040&240000&48640\\
1&582160&1108000&1118304&557584&110864\\
2&690608&1195520&1018512&418768&68128\\
3&756896&1304520&1151162&522393&96703\\
4&516296&613195&462027&220948&46208\\
5&181031&120245&477&-23717&-6797\\
6&-259&2795&1978&774&129\\
7&43&0&0&0&0
% HD-CURVATURE-END
\end{array}
\]
This displayed identity proves the sign without a parameter certificate:
$P_0,\ldots,P_4,P_7$ have positive coefficients, and, for $0\le y<1$,
\[
 P_5(y)\ge181031-23717-6797=150517,\qquad P_6(y)\ge-259.
\]
Thus $P_5+RP_6\ge150517-8\cdot259>0$. Hence $Q>0$ and $J$ strictly decreases.
As $W''=(2\gamma+6\beta/s)(J/L-1)$, the asserted curvature pattern follows.

To turn this curvature pattern into a supporting line, subtract the
tangent at $p_*$ and write
$D(p)=W(p)-W(p_*)-W'(p_*)(p-p_*)$.
Then $D(p_*)=D'(p_*)=0$, and $D'$ first increases and then decreases.
If $D(1/4)<0$, the contact $p_*$ must lie in the strictly concave portion:
otherwise $D'<0$ throughout its left side, which would give
$D(1/4)>0$. In the concave portion, the tangent strictly
dominates to the right of $p_*$; to its left, $D$ has no maximum except
at $1/4$ or $p_*$. Thus $D<0$ away from $p_*$.
In the present case the
left-endpoint gap is strictly positive, so $D(1/4)<0$:
\[
 \Delta=-D(1/4)=h_q+\frac k8-2+\frac\beta{12}+\frac\gamma{144}
 =h_q+\frac k8-2+\frac{13}{432}>\frac1{5000}.
 \tag{HUC}\label{eq:fqe-hs-uniform-contact}
\]
For completeness, this small constant has an elementary rational proof.
Let $u_0=19331/10000$, $v=1/(2u_0+1)$ and
\[
 L_-=2\sum_{j=0}^5\frac{1}{(2j+1)3^{2j+1}},\qquad
 B_+=2\sum_{j=0}^4\frac{v^{2j+1}}{2j+1}
             +\frac{2v^{11}}{11(1-v^2)}.
\]
The formula $h_q=8/3-(5/(3L))\ln(1+1/u)$ and the atanh series give $L>L_-$ and
$\ln(1+1/u_0)<B_+$. Since $u_0^5<27$, $u>u_0$; the function
$-5\ln(1+1/u)+2(u-1)/(u+1)$ increases with $u>0$.
The purely rational comparison
\[
 \frac{301}{144}L_- -5B_+ +\frac{2(u_0-1)}{u_0+1}
 >\frac{21}{50000}
\]
therefore gives $3L\Delta>21/50000>3L/5000$, using $L<7/10$.
This proves~\eqref{eq:fqe-hs-uniform-contact}.
The tangent bound is strict away from $p_*$ by the strict curvature
calculation. Thus $\Lambda\le h_q$, with equality only at HS.

Finally $u>19/10$ and $L<7/10$ give
$k>480/203>7/3$. Substitution into the outer dual yields
\[
 \ell=\left(\alpha-\gamma(\tfrac{17}{12}+\tfrac\varepsilon3)\right)\varepsilon
          +F\left(\tfrac1{12}-\tfrac59F\right)
 \ge\frac{149}{324}\varepsilon+F\left(\tfrac1{12}-\tfrac59F\right).
\]
The last term is nonnegative for $F\le3/20$, as required.
All global maximisation and sign arguments above are analytic;
HD checks the displayed identities and rational comparisons.
\end{proof}


\subsection{The strict supporting band}\label{sec:fqe-rational-dual}
The bounded-head and capped-cubic proofs below use the same near-uniform
moment inequality. Readers following dependencies rather than the
case partition can first consult
Lemma~\ref{lem:fqe-near-uniform-moment} and the moment-rectangle argument
in Section~\ref{sec:sup-nu-moment}, then return to these two applications.
\begin{lemma}[Strict dual band; analytic with exact polynomial sign certificates]
\label{lem:fqe-rational-dual}
For band $L$ of Table~\ref{tab:fqe-dual-co}, set
$W=5/3+\min\ell_\theta/3$, with the minimum over its rectangle's corners.
At $q_s=1587/1000$, every four-point spectrum with $p\le49/75$ satisfies
$\Lambda_\theta<W-10^{-5}$.
\end{lemma}
\begin{proof}
At fixed purity the variable part is $H_q-(9/5)t$.
Lemma~\ref{lem:fqe-low-tilt-reduction} reduces the maximum to $(a,b,c,c)$.
Lemma~\ref{lem:fqe-small-head-dual} handles $a\le1/2$ and
Lemma~\ref{lem:fqe-large-head-dual} handles $1/2\le a\le4/5$.
These cover the domain: $p\ge a^2+(1-a)^2/3$ and its value $49/75$ at $a=4/5$
force $a\le4/5$. The corner minimum of the outer dual is $1837/125000$,
so $W-10^{-5}=2507333/1500000$ as required.
\end{proof}

\begin{lemma}[Binary quartic correction; analytic]\label{lem:fqe-binary-quartic}
Let $1<q<2$. Put $v=(1-2u)^2$,
\[
 m_q(v)=\frac{(1+\sqrt v)^q+(1-\sqrt v)^q}{2},\qquad
 D_q(v)=\frac{\ln m_q(v)}{(q-1)\ln2}=1-H_q(u,1-u).
\]
Then $D_q'''(v)>0$ on $0<v<1$, and
\[
 A_q(u)=\frac{4u(1-u)-H_q(u,1-u)}{u(1-u)(1-2u)^2}
\]
is strictly increasing on $0<u<1/2$.
Consequently, a quartic bound
$H_q(u,1-u)\le4u(1-u)-a u(1-u)(1-2u)^2$
at $u=u_0$ holds throughout $u_0\le u\le1/2$.
\end{lemma}
\begin{proof}
Write $a=\tan^2\!\bigl(\tfrac{\pi}{2q}\bigr)$; since $\pi/4<\pi/(2q)<\pi/2$, $a>1$. Put $f=\ln m_q$.
The proof is Cauchy's formula for $\ln m_q$ on the slit plane, after removing the one zero of $m_q$.

\emph{Analyticity.} Let $g(w)=\tfrac12\bigl((1+w)^q+(1-w)^q\bigr)$ with principal powers. It is holomorphic on
$\mathbb C\setminus\bigl((-\infty,-1]\cup[1,\infty)\bigr)$ and even, so $m_q(v)=g(\sqrt v)$ does not depend on the
sign of the root and is holomorphic on $\Omega=\mathbb C\setminus[1,\infty)$; at $v=0$ it is the convergent series
$\sum_k\binom q{2k}v^k$. It is real on the real axis and $m_q(\bar v)=\overline{m_q(v)}$.
For $\operatorname{Im}v\ge0$ take $w=\sqrt v$ in the closed first quadrant. Then $\arg(1-w)=\arg(w-1)-\pi$, so
\begin{equation}
 m_q(v)=\tfrac12\bigl[(1+w)^q+e^{-i\pi q}(w-1)^q\bigr],\qquad \arg(w-1)\in[0,\pi].
 \label{eq:fqe-bq-upper}
\end{equation}
The right side is continuous on the closed first quadrant, so $m_q$ extends continuously to the closed upper
half-plane, and to the closed lower half-plane by conjugation.

\emph{Zeros.} Let $m_q(v)=0$ and $\operatorname{Im}v\ge0$, $w=\sqrt v$. Then $|1+w|^q=|1-w|^q$, so $\operatorname{Re}w=0$ and $v=-y^2$
is real, $w=iy$, $y\ge0$. Now $(1\pm iy)^q=(1+y^2)^{q/2}e^{\pm iq\arctan y}$, so
$m_q(-y^2)=(1+y^2)^{q/2}\cos(q\arctan y)$. As $0<q\arctan y<q\pi/2<\pi$ for $y>0$, this vanishes only at
$q\arctan y=\pi/2$, that is $v=-a$. (For $v\in[0,1)$, $m_q>0$.) The zero is simple:
$g'(w)=\tfrac q2\bigl((1+w)^{q-1}-(1-w)^{q-1}\bigr)$, and at $w=iy_0$, $y_0=\sqrt a$,
\[
 g'(iy_0)=iq(1+y_0^2)^{(q-1)/2}\sin\bigl((q-1)\arctan y_0\bigr)\ne0,
\]
since $0<(q-1)\arctan y_0<\pi/2$; and $m_q'(v)=g'(w)/(2w)$. Hence
$h(v)=m_q(v)/(1+v/a)$ is holomorphic and zero-free on $\Omega$, and continuous and zero-free on the closed upper
half-plane away from the cut; on the cut see the next step.

\emph{Boundary values on the cut.} For $t>1$, \eqref{eq:fqe-bq-upper} gives
$2m_q(t+i0)=(1+\sqrt t)^q+e^{-i\pi q}(\sqrt t-1)^q$ with
$\operatorname{Im}\,2m_q(t+i0)=-\sin(\pi q)(\sqrt t-1)^q>0$, because $\pi<\pi q<2\pi$.
So $m_q(t+i0)\ne0$ and
\[
 \theta_q(t)=\arg\bigl(2m_q(t+i0)\bigr)\in(0,\pi)
\]
is continuous on $[1,\infty)$, with $\theta_q(t)\to0$ as $t\downarrow1$, since $m_q(1)=2^{q-1}>0$.

\emph{The logarithm.} $\Omega$ is simply connected and $h$ is zero-free on it with $h>0$ on $(-\infty,1)$
(there $m_q$ and $1+v/a$ change sign together at $-a$). So $F=\ln h$ with $F(0)=0$ is holomorphic on $\Omega$,
real on $(-\infty,1)$, and $F(\bar v)=\overline{F(v)}$. The closed upper half-plane is simply connected and $h$
is continuous and zero-free on it, so $F$ extends continuously to it. On the cut, $\operatorname{Im}F(t+i0)$ is
continuous and equals $\theta_q(t)$ modulo $2\pi$ (as $1+t/a>0$). Both vanish at $t=1$, so
$\operatorname{Im}F(t+i0)=\theta_q(t)$ for all $t\ge1$ and $F(t+i0)-F(t-i0)=2i\theta_q(t)$.
In particular $F$ is bounded near $v=1$.

\emph{Growth.} For $|w|>1$ in the first quadrant, $(1+w)^q=w^q(1+1/w)^q$ and $(w-1)^q=w^q(1-1/w)^q$ (the
arguments add without wrapping), so by \eqref{eq:fqe-bq-upper}
\[
 m_q(v)=w^q\bigl(c+\varepsilon(1/w)\bigr),\qquad c=\tfrac12\bigl(1+e^{-i\pi q}\bigr)\ne0,
\]
with $\varepsilon$ holomorphic on $|z|<1$ and $\varepsilon(0)=0$; $c\ne0$ because $q$ is not an odd integer.
Hence $F(v)-\bigl(\tfrac q2\mathrm{Log}\,v-\mathrm{Log}(1+v/a)+\mathrm{Log}\,c\bigr)$ has exponential tending to
$1$ as $|v|\to\infty$ in the closed half-plane, which is connected, so it tends to a constant in $2\pi i\mathbb Z$.
Thus $|F(v)|\le C\ln|v|$ for large $|v|$, in both half-planes.

\emph{The integral formula.} Fix $0<v<1$. By Cauchy's formula for $F'''$ on the region
$\{|\zeta|<R\}$ with the strip $\{\operatorname{Re}\zeta\ge1,\ |\operatorname{Im}\zeta|\le\epsilon\}$ removed,
\[
 F'''(v)=\frac{3!}{2\pi i}\oint\frac{F(\zeta)}{(\zeta-v)^4}\,d\zeta .
\]
Let $\epsilon\downarrow0$ and then $R\to\infty$. $F$ is bounded near $1$, so the short segment near the
tip $\zeta=1$ contributes $O(\epsilon)$; the arc contributes $O(R^{-3}\ln R)$; and by the jump of $F$ across the cut
the two edges contribute $2i\int_1^R\theta_q(t)(t-v)^{-4}\,dt$. Therefore
\[
 F'''(v)=\frac6\pi\int_1^\infty\frac{\theta_q(t)}{(t-v)^4}\,dt>0 .
\]
Since $m_q=(1+v/a)e^F$ and $(\ln(1+v/a))'''=2/(a+v)^3$,
\[
 f'''(v)=(\ln m_q)'''(v)=\frac2{(a+v)^3}+\frac6\pi\int_1^\infty\frac{\theta_q(t)}{(t-v)^4}\,dt>0
 \qquad(-a<v<1),
\]
which is $D_q'''>0$ on $0<v<1$.

\emph{Monotonicity of $A_q$.} $D_q(0)=0$ and $D_q(1)=1$, so
\[
 \frac{A_q(u)}4=\frac{D_q(v)-v}{v(1-v)}=-[0,v,1]D_q .
\]
For $0<v<v'<1$, $[0,v,1]D_q-[0,v',1]D_q=(v-v')[0,v,v',1]D_q$, and by the mean value theorem for
divided differences (Rolle's theorem; $D_q$ is continuous on $[0,1]$) the last factor is
$D_q'''(\xi)/6>0$. So $A_q$ decreases in $v=(1-2u)^2$: it is strictly increasing in $u$.
\end{proof}


\subsection{The bounded-head chart}
\label{sec:fqe-small-head-dual}

\begin{lemma}[Supporting bound on the bounded-head chart; analytic; exact rational checks]\label{lem:fqe-small-head-dual}
At $q=1587/1000$ and the multiplier $\theta=(397/100,3/20,159/100)$,
every ordered spectrum $(a,b,c,c)$ with $a\le1/2$ satisfies
\[
 \Lambda_\theta\le\frac{8023}{4800}
 <\frac53+\frac{1837}{375000}-10^{-5}.
\]
So on this chart the band-$L$ dual bound needs only rational arithmetic: rational enclosures of eight node values (SHD-1) for the binary
quartic correction of Lemma~\ref{lem:fqe-binary-quartic} and eight one-variable quartic sign checks (SHD-3).
\end{lemma}
\begin{proof}
Put $r=q-1$, $L=\ln2$, $h(u)=H_q(u,1-u)$, $d=u(1-u)$, and write
\[
 \lambda=\Bigl(\frac{1-u+\delta}2,\frac{1-u-\delta}2,\frac u2,\frac u2\Bigr),\qquad v=\delta^2,
\]
a head pair of total weight $1-u$ and a uniform tail pair of weight $u$.
The head cap and ordering give $0\le u\le1/2$ and $0\le\delta\le\min\{u,1-2u\}$.
We treat $u\le1/4$ by a head-deficit estimate and $u\ge1/4$ by the near-uniform moment bound.

\emph{Small tail, $0\le u\le1/4$.}
For the conditional vectors of the two pairs, Jensen's inequality for the logarithm (as in the proof
of~\eqref{eq:fqe-capped-pairs}) gives
\[
 H_q(\lambda)\le h(u)+(1-w)\bigl(1-D_q(z)\bigr)+w
 =1+h(u)-(1-w)D_q(z),\qquad w=\frac{u^q}{(1-u)^q+u^q},
\]
where $z=v/(1-u)^2\le u^2/(1-u)^2\le1/9$ is the squared imbalance of the head pair and $D_q$ is the
function of Lemma~\ref{lem:fqe-binary-quartic}. That lemma shows that $(D_q(z)-z)/(z(1-z))=A_q/4$ decreases in
$z\in(0,1)$. Put $\kappa=\tfrac{81}8\bigl(D_q(1/9)-\tfrac19\bigr)$, its value at $z=1/9$. Here
$D_q(1/9)=1-H_q(1/3,2/3)\ge249/2000>1/9$, since $H_q(1/3,2/3)<1751/2000$ (power sums and the atanh series, as for the
nodes below), so $\kappa>0$. For $0<z\le1/9$,
\[
 D_q(z)\ge z+\kappa z(1-z)\ge z\bigl(1+\tfrac89\kappa\bigr)=9D_q(1/9)\,z\ge\frac{2241}{2000}z ,
\]
since $9\cdot249/2000=2241/2000$.

Next $1-w=(1-u)^2k_q(u)$ with
\[
 k_q(u)=\frac{(1-u)^{q-2}}{(1-u)^q+u^q}=\frac{(1-u)^{-2}}{1+(u/(1-u))^q},
\]
which increases with $q$ because $u\le1-u$; hence $k_q\ge k_{3/2}$. With $t=\sqrt{u/(1-u)}$,
\[
 \frac{k_{3/2}(u)}{1+2u}=\Bigl[1+\Bigl(\frac{t(1-t)}{1+t^2}\Bigr)^3\Bigr]^{-1}\ge\frac{64}{65},
\]
since $0\le t(1-t)\le1/4$ and $1+t^2\ge1$, so the cube is at most $1/64$.
Combining, $(1-w)D_q(z)\ge\frac{2241}{2000}\cdot\frac{64}{65}(1+2u)v=\frac{8964}{8125}(1+2u)v$, so
\[
 H_q(\lambda)\le1+h(u)-\frac{11}{10}(1+2u)v.                 \tag{SHD}\label{eq:fqe-small-head-deficit}
\]

The binary entropy is controlled by the quartic correction of Lemma~\ref{lem:fqe-binary-quartic}:
$h(u)=4d-A_q(u)\,d(1-4d)$ with $1-4d=(1-2u)^2\ge0$, and $A_q$ is increasing. Take the nodes and values
\[
\begin{array}{c|cccccccc}
u_i&0&\frac1{100}&\frac3{100}&\frac1{20}&\frac2{25}&\frac{13}{100}&\frac{39}{200}&\frac{47}{200}\\[2pt]\hline
a_i&\frac{99}{1000}&\frac{217}{1000}&\frac{289}{1000}&\frac{83}{250}&\frac{377}{1000}&\frac{43}{100}&\frac{12}{25}&\frac{101}{200}
\end{array}
\]
and $u_8=1/4$. They satisfy $A_q(u_i)\ge a_i$: at $u_0=0$ in the sense $A_q(0^+)=4-q/(rL)>99/1000$, and at the other seven
points by interval enclosures of the power sum $u^q+(1-u)^q$ and of the logarithm (atanh series) [exact certificate; SHD-1].
Since $A_q$ is increasing, $A_q\ge a_i$ on $[u_i,u_{i+1}]$, hence $h(u)\le4d-a_id(1-4d)$ there.

Substituting~\eqref{eq:fqe-small-head-deficit} in the definition of $\Lambda_\theta$ gives
\[
 \Lambda_\theta\le\frac{1999}{1200}+h(u)-\frac{179}{50}d
       -\frac{159}{100}d^2
       +\left(\frac3{200}+\frac{37}{50}u-\frac{159}{100}u^2\right)v
       -\frac{159}{400}v^2.
\]
The coefficient of $v$ is positive on $[0,1/4]$. Using $v\le u^2$, dropping the last term and inserting the bound for $h$ leaves, on $[u_i,u_{i+1}]$, the quartic
\[
 P_i(u)=\frac{1999}{1200}+\frac{21}{50}d-a_i\,d(1-4d)-\frac{159}{100}d^2
       +\Bigl(\frac3{200}+\frac{37}{50}u-\frac{159}{100}u^2\Bigr)u^2 .
\]
It remains to show $P_i\le8023/4800$ on $[u_i,u_{i+1}]$. Each $8023/4800-P_i$ is a quartic in $u$; written in the
Bernstein basis of its interval, all five coefficients are positive [exact-rational certificate; SHD-3]. The smallest coefficient is $1/6400$, at $u=1/4$.

\emph{Large tail, $1/4\le u\le1/2$.} Here we show that $\Lambda_\theta$ does not exceed its uniform value.
The coordinate bounds give $p\le11/32$, by maximising a convex
quadratic separately on $[1/4,1/3]$ and $[1/3,1/2]$.
Put $x=p-1/4$, $y=\sum(\lambda_i-1/4)^3$.
Here $y\ge0$, so $v=y/x$ lies in $[0,9/50]$, and $x\le X=3/32$.
With $b=12\beta=9/5$ and $c=\alpha+8\beta-\gamma/2=35/8$, the identity~\eqref{eq:fqe-uniform-dual-difference},
valid for every multiplier, gives $\Lambda_\theta((1/4)^4)-\Lambda_\theta(\lambda)=2-H_q-cx+by+\gamma x^2$.
We apply Lemma~\ref{lem:fqe-moment-rectangle} with the SHD column of Table~\ref{tab:fqe-moment-rectangle}.
On the stated rectangle,
$17/10<V<19/10$, $-3/50<W<3/8$, $1\le N<6/5$ and $|B+2Cv+XD_5|<19/10$,
with $A,B,C,D_5$ those of Lemma~\ref{lem:fqe-near-uniform-moment}.
The bound $V^2-2W>8/3$ from the capped-cubic proof gives $Q=8/3$, and
$8/3-2X(19/10)(3/50)>(7/4)(6/5)^2$ gives $m=7/4>2\rho\gamma$; $C>7/5$ gives
$\Phi_{vv}(X,v)\ge(2C-X(19/10)^2)/(\rho(6/5)^2)>4=\kappa$.
At the single point $v_*=27/200$,
\[
 \Phi(X,v_*)>1/800,\qquad |\Phi_v(X,v_*)|<1/100,
\]
using the eight-term alternating lower sum for
$\ln(1+XV+X^2W)$ and $693147/10^6<L<693148/10^6$.
The lemma gives a dual difference greater than $(1/800-1/80000)x>0$.
Thus $\Lambda_\theta\le\Lambda_\theta((1/4)^4)=8023/4800$.
The uniform case $x=0$ is immediate.

The small-tail and large-tail cases together prove the stated bound.
\end{proof}

\subsection{The large-head chart by two Hermite bounds}
\label{sec:fqe-large-head-dual}
Throughout this subsection $q=1587/1000$, $r=q-1$, $L=\ln2$ and
$\theta=(397/100,3/20,159/100)$. Put
\begin{gather*}
 f(a,z)=\Lambda_\theta\left(a,\frac{(1-a)(1+2z)}3,
              \frac{(1-a)(1-z)}3,\frac{(1-a)(1-z)}3\right),\\
 B=\frac{1+2z^2}3,\qquad C=\frac{1+6z^2+2z^3}9,\qquad
 K=\frac{(1+2z)^q+2(1-z)^q}{3^q}.
\end{gather*}
Here $1/2\le a\le4/5$, $0\le z\le1$, and the purity and third moment are
$p=a^2+B(1-a)^2$, $t=a^3+C(1-a)^3$.
The objective is $f<T=2507333/1500000$. The first lemma supplies the
shape condition $f_{aaa}\le0$. The second then bounds $f$ by a quadratic
in $a$, replaces $K(z)$ by one cubic minorant, and proves positivity of
that quadratic's gap by completing the square. Thus the two-dimensional
chart reduces to univariate polynomial signs.

\begin{lemma}[Largest-entry curvature; exact polynomial sign certificates]\label{lem:fqe-large-head-curvature}
On this whole chart, $f_{aaa}\le0$.
\end{lemma}
\begin{proof}
Write $h(a,K)=-\ln(a^q+K(1-a)^q)/(rL)$. The remaining part of $f$ is a polynomial in $a$
whose only fourth derivative comes from $-\gamma p^2$, so
$f_{aaaa}=h_{aaaa}-24\gamma(1+B)^2$.

\emph{The shape of $h_{aaaa}$ in $K$.}
Set $x=2a-1\in[0,3/5]$, $\rho=(1-a)/a=(1-x)/(1+x)\in[1/4,1]$, $w=K\rho^q$ and
\[
 y=\frac{a^q-K(1-a)^q}{a^q+K(1-a)^q}=\frac{1-w}{1+w}.
\]
Apart from the positive factor $16q/(rL(1-x^2)^4)$, $h_{aaaa}$ is
\begin{align*}
 J(y)={}&6q^3y^4-24q^2xy^3+(-8q^3+36qx^2+8q)y^2\\
 &+(24q^2x-24x^3-24x)y+2q^3-36qx^2-8q+6x^4+36x^2+6.
\end{align*}
On $0\le x\le3/5$,
\[
 J'(y)\big|_{y=x}=8x(q-1)\{3(q-1)^2x^2-2q^2+q+3\}\le0,\qquad
 J'''(y)=144q^2(qy-x)\ge0\quad(y\ge x).
\]
So on $[x,1]$ the derivative $J'$ is convex and starts nonpositive. Hence $J$ decreases and then
increases, and has no interior maximum on any subinterval:
\[
 J(y)\le\max\{J(y_\ell),J(y_h)\}\qquad(x\le y_\ell\le y\le y_h\le1).
\]

\emph{Envelopes for $y$.}
Write $K=\lambda_1^q+2\lambda_2^q$ with $\lambda=\bigl(\frac{1+2z}3,\frac{1-z}3,\frac{1-z}3\bigr)$.
Convexity of $s^q$ gives $K\ge3^{-r}>1/2$, since $3^r<2$.
Concavity of $s^r$ with weights $\lambda_i$ gives $K=\sum\lambda_i\lambda_i^r\le B^r\le1$.
If $B\le1/2$, then $K\le2^{-r}<2/3$, because $2^q>3$.
For $0<\rho\le1$ we have $\rho^2\le\rho^q\le\rho^{3/2}\le\rho(1+\rho)/2$ and $\rho^q\le\rho$.
Hence $w\ge\rho^2/2$, $w\le\rho(1+\rho)/2=(1-x)/(1+x)^2$ always, and $w\le2\rho/3$ when $B\le1/2$.
Since $y=(1-w)/(1+w)$ decreases in $w$, we get $y_\ell\le y\le y_h$ with
\[
 y_h=\frac{(1+x)^2-(1-x)^2/2}{(1+x)^2+(1-x)^2/2},\quad
 y_A=\frac{(1+x)-\frac23(1-x)}{(1+x)+\frac23(1-x)},\quad
 y_B=\frac{(1+x)^2-(1-x)}{(1+x)^2+(1-x)},
\]
and $y_\ell=y_A$ if $B\le1/2$, $y_\ell=y_B$ otherwise.
Moreover $y_A-x$ and $y_B-x$ have the numerators $\frac13(1-x^2)$ and $(1-x)(x+x^2)$ over the
positive denominators, so $x\le y_\ell$, and $y_h\le1$.

\emph{Three polynomial inequalities.}
Let $L_0=0.693147<L$. For $0\le x\le3/5$,
\begin{align*}
 J(x,y_A)&\le\tfrac{1696}{25}\,\frac{rL_0(1-x^2)^4}{16q},\qquad
 J(x,y_h)\le\tfrac{1696}{25}\,\frac{rL_0(1-x^2)^4}{16q},\\
 J(x,y_B)&\le\tfrac{4293}{50}\,\frac{rL_0(1-x^2)^4}{16q}.
\end{align*}
Clearing the denominators ($D_A^4$ with $D_A$ linear, $D_B^4$ and $D_h^4$ with quadratic denominators)
turns each into a polynomial inequality of degree $12$, $16$ and $16$ in $x$.
Each is certified by the Bernstein expansion on $[0,3/5]$, all of whose coefficients are positive (LHD-1).
Thus $h_{aaaa}<1696/25=24\gamma(16/9)$ if $B\le1/2$, and $h_{aaaa}<4293/50=24\gamma(9/4)$ otherwise.
Since $B\ge1/3$, respectively $B\ge1/2$, this is $h_{aaaa}<24\gamma(1+B)^2$, that is, $f_{aaaa}<0$.

\emph{The point $a=1/2$.}
There, with $v=(1-K)/(1+K)$,
\[
 f_{aaa}=\frac{16q(q+1)v-16q^3v^3/r}{L}
             -\frac{54}5(1-C)-\frac{477}{25}(1-B^2).
\]
The first term increases with $v\le1/3$. Since $K^{(4)}>0$ (because $q(q-1)(q-2)(q-3)>0$),
cubic Hermite interpolation at $0,1$ gives, with $k=5247/10000<K(0)=3^{-r}$,
\[
 K\ge K_*(z)=k+(3-3k-2q/3)z^2+(2q/3-2+2k)z^3.
\]
Replacing $K$ by $K_*$ and $L$ by $L_0$, and clearing
$r(1+K_*)^3$, reduces $-f_{aaa}(1/2,z)\ge0$ to $(1-z)\sum_{j=0}^{12}c_jz^j\ge0$, where every $c_j$ is
positive (LHD-2). Lowering $K(0)$ to $k$ preserves the minorant because its coefficient is
$1-3z^2+2z^3=(1-z)^2(1+2z)\ge0$. Together with $f_{aaaa}<0$ this proves the lemma.
\end{proof}

\begin{lemma}[The whole large-head chart; exact polynomial sign certificates]\label{lem:fqe-large-head-dual}
Every spectrum on the above chart satisfies
$\Lambda_\theta<T=2507333/1500000$.
\end{lemma}
\begin{proof}
Put $a_0=127/200$, $S=1-a_0=73/200$, $d=a_0-1/2=27/200$ and $\xi=(a-a_0)/d\in[-1,11/9]$.
By Lemma~\ref{lem:fqe-large-head-curvature}, the quadratic Hermite interpolant in $a$
at the single node $1/2$ and the double node $a_0$ is a global upper bound,
since its remainder has the sign of $f_{aaa}(\eta,z)(a-1/2)(a-a_0)^2\le0$:
\[
 f(a,z)\le Q=f_0(1-\xi^2)+f_l\xi^2+dD(\xi+\xi^2),\quad
 f_0=f(a_0,z),\ f_l=f(\tfrac12,z),\ D=f_a(a_0,z).
\]
Write $f=h(a,K)+P$, where $P$ is a polynomial in $(a,B,C)$ and does not depend on $K$.
Then only the three values $h(a_0,K)$, $h(\tfrac12,K)$, $h_a(a_0,K)$ in $Q$ depend on $K$.

\emph{$Q$ decreases in $K$ on $[.52,1]$.}
Let $R=(a_0/S)^q\in[12/5,5/2]$, $c=rL$ and $\lambda=dq/(a_0S)=42849/46355$.
Differentiating,
\[
 h_K(a_0,K)=\frac{-1}{c(R+K)},\quad h_K(\tfrac12,K)=\frac{-1}{c(1+K)},\quad
 h_{aK}(a_0,K)=\frac{qR}{c\,a_0S(R+K)^2},
\]
and therefore
\[
 -c(R+K)^2(1+K)\,Q_K=(1+K)\,(R+K-\lambda R\xi)+\xi^2\{(R+K)(R-1)-\lambda R(1+K)\}.
\]
For $-1\le\xi\le11/9$ both terms are nonnegative, the first strictly positive. The first is linear in $\xi$, positive at both ends:
$\frac{11}9\lambda<1.13$ and $K>.13R$. The brace increases in $R$ and decreases in $K$, and is
$\frac{17}5\cdot\frac75-\frac{24}5\lambda>0$ at $R=12/5$, $K=1$ (LHD-3).

\emph{A cubic minorant of $K$.}
Because $K^{(4)}>0$, cubic Hermite interpolation at the double nodes $3/20,17/20$ gives
\[
 K(z)\ge\kappa(z)=.5238033+.0142244z+.4204281z^2+.0384307z^3,
 \qquad .52<\kappa<1.                                      \tag{LH1}
\]
Each coefficient of $\kappa$ is a rational lower rounding of the interpolant's coefficient (LHD-3).
Since $.52<\kappa\le K\le1$, we have $Q(K)\le Q(\kappa)$.

\emph{Rational bounds for the logarithms.} Set
\[
 \rho=2.40792594,\quad u=2.45774282,\quad H_0=1.10489073,\quad H_l=1.32818582.
\]
They approximate $R$, $1/(rL)$, $h(a_0,3/4)$, $h(\tfrac12,3/4)$ respectively with error less than $2\cdot10^{-8}$.
For $v\in[.52,1]$ define
\[
 \widehat h_0=H_0-2u\frac{v-3/4}{2\rho+v+3/4},\quad
 \widehat h_l=H_l-2u\frac{v-3/4}{v+11/4},\quad
 \widehat h_a=-\frac{2qu}{1-(27/100)^2}\Bigl(\frac{\rho-v}{\rho+v}-\frac{27}{100}\Bigr).
\]
The identity $\ln((1+y)/(1-y))=2y+e_y$ has $|e_y|\le2|y|^3/(3(1-y^2))$.
Here $h(a_0,v)-h(a_0,\tfrac34)=-\frac1{rL}\ln\frac{R+v}{R+3/4}$ and
$h(\tfrac12,v)-h(\tfrac12,\tfrac34)=-\frac1{rL}\ln\frac{1+v}{7/4}$. The logarithms equal $\ln\frac{1+y}{1-y}$ with
$y=(v-\frac34)/(2R+v+\frac34)$, respectively $y=(v-\frac34)/(v+\frac{11}4)$, and $|y|<.0383$, respectively $|y|<.0704$.
Rounding the constants and the remainder give
\[
 |h(a_0,v)-\widehat h_0|<e_0=10^{-4},\quad
 |h(\tfrac12,v)-\widehat h_l|<e_l=6\cdot10^{-4},\quad
 |h_a(a_0,v)-\widehat h_a|<e_a=10^{-6}.                         \tag{LH2}
\]
These include the rounding errors; LHD-3 checks the scalar arithmetic.
On $[-1,11/9]$ we have $|1-\xi^2|\le1$ and $|d(\xi+\xi^2)|\le11/30<1$, so with
$P_0=P(a_0)$, $P_l=P(\tfrac12)$ and $P_a=\partial_aP(a_0)$,
\[
 Q(\kappa)\le F=(\widehat h_0+P_0)(1-\xi^2)+(\widehat h_l+P_l)\xi^2+d(\widehat h_a+P_a)(\xi+\xi^2)
   +e_0+e_a+e_l\xi^2,\qquad v=\kappa(z).
\]

\emph{The completed square.}
Put $E=(v+2\rho+3/4)(v+11/4)(v+\rho)>0$. Then $E\,(T-F)=c_0+c_1\xi+c_2\xi^2$, where after
substituting $v=\kappa(z)$ and the polynomials $B,C$ the $c_i$ are polynomials in $z$.
On $[0,1]$,
\[
 c_2\ge\tfrac1{10},\qquad 4c_0c_2-c_1^2\ge\tfrac3{100},
\]
because all Bernstein coefficients of $c_2$ (degree $13$) and of $4c_0c_2-c_1^2$ (degree $26$) are at least
these values (LHD-4). Thus the quadratic in $\xi$ is positive for every real $\xi$, so $f\le Q\le F<T$.
\end{proof}


\subsection{A strengthened cubic bound below purity $47/100$}
\label{sec:fqe-capped-cubic}
We now fix $q=1587/1000$, $r=q-1$, $L=\ln2$ and $x=p-1/4$.
\begin{lemma}[Capped cubic slack; analytic with explicit scalar arithmetic]
\label{lem:fqe-capped-cubic}
Every four-point probability vector with $p\le47/100$ satisfies
\[
H_q(\lambda)\le \frac{10}3-6p+\frac83t
-\frac{11}{10}(p-1/4)(1/2-p).
\tag{CG}\label{eq:fqe-capped-cubic}
\]
Equality holds only at the uniform-four spectrum.
\end{lemma}
\begin{proof}
In centred moments $y=\sum_i(\lambda_i-1/4)^3$, the assertion is
\[
2-H_q\ge \frac{171}{40}x-\frac83y-\frac{11}{10}x^2.
\]
We first prove it for $p\le21/50$, then use binary grouping for the rest.

\emph{The moment argument, $x\le17/100$.}
The claim follows from~\eqref{eq:fqe-moment-rectangle} of Lemma~\ref{lem:fqe-moment-rectangle} with $b=8/3$,
$c=171/40$, $\gamma=11/10$ and $\mu>0$; we apply that lemma with the CG column of
Table~\ref{tab:fqe-moment-rectangle}. Here $N=1+xV+x^2W$, $V=A+Bv+Cv^2$, $W=C/4+D_5v$, $v=y/x$, and
$A,B,C,D_5$ are those of Lemma~\ref{lem:fqe-near-uniform-moment}; the rectangle lemma does not require
its smaller purity domain. The third-moment bound gives $|v|<6/25$.
On $0\le x\le X=17/100$, $|v|\le6/25$, direct rational bounds give
\[
5/3<V<9/4,\quad -21/100<W<1,\quad 1\le N<3/2,\qquad
V^2-2W\ge\frac{1706}{625}-\frac9{10}|v|+5v^2>\frac83,
\]
and $|B+2Cv+XD_5|<213/100$. Then $(8/3-2X(9/4)(21/100))/(3/2)^2>1=m$, so $(\ln N)_{xx}<-1$, and
$m>(11/5)\rho=2\rho\gamma$; $C>144/100$ gives $\Phi_{vv}(X,v)\ge(2C-X(213/100)^2)/(\rho(3/2)^2)>2=\kappa$.
At $v_*=0$ put
$z=XA+X^2C/4$, use the eight-term lower sum
$\ln(1+z)\ge\sum_{j=1}^8(-1)^{j+1}z^j/j$, and use
$693147/10^6<L<693148/10^6$: then $\Phi(X,0)>1/250$ and $|\Phi_v(X,0)|<1/20$.
The lemma gives $\mu\ge1/250-(1/20)^2/4>0$. This proves the claim strictly
when $x>0$; $x=0$ is the uniform spectrum.

\emph{Two binary bounds from endpoint values.}
Write $h(u)=H_q(u,1-u)$ and $d_0=u(1-u)$.
The monotonicity of $A_q$ in Lemma~\ref{lem:fqe-binary-quartic} (proved there by Cauchy's formula) gives
\[
 h(u)\le4d_0-\frac{11}{40}d_0(1-4d_0)\quad(3/100\le u\le1/2),
\qquad
 h(u)\le4d_0-\frac38d_0(1-4d_0)\quad(1/12\le u\le1/2).
\tag{BQ}\label{eq:fqe-capped-binary}
\]
Only the two lower endpoints require checking. Explicit sufficient bounds are
\[
\begin{array}{c|c|c|c}
u_0&u_0^q>&(1-u_0)^q>&h(u_0)<\\\hline
3/100&19/5000&1191/1250&1093/10000\\
1/12&193/10000&871/1000&5713/20000
\end{array}
\]
The power bounds follow by raising to the thousandth power. The power sums are then at least
$S=4783/5000$ and $S=8903/10000$, and the entropy caps follow from
$-\ln(1-z)\le z+z^2/2+z^3/(3(1-z))$ with $z=1-S$; the two quartics exceed them (CG-2).

It follows that $4u(1-u)-h(u)>(11/10)d(1/4-d)$ whenever
$3/100\le u\le4/15$ and $3/100\le d\le\min\{2/25,d_0\}$.
For $u\le1/12$, use the first quartic and the fact that
$d(1/4-d)$ increases on $[0,1/8]$.
For $u\ge1/12$, use the second quartic: the concave function $d_0(1-4d_0)$
has its minimum at an endpoint of $11/144\le d_0\le44/225$, and
\[
\frac38\min\left\{\frac{11}{144}\frac{25}{36},
\frac{44}{225}\frac{49}{225}\right\}
>\frac{187}{12500}
\ge\frac{11}{10}d(1/4-d).
\]

\emph{The two pairs, $21/50\le p\le47/100$.}
Determinant monotonicity preserves $p,t$ and increases $H_q$, so it suffices
to use the two ordered charts in \eqref{eq:fqe-sliver-coordinates}.
Put $v=\delta^2$, $d=1/2-p=u-u^2-v/2\in[3/100,2/25]$.
Ordering gives $3/100<u<4/15$ on S and $3/100<u<3/32$ on B.
For S use $p\le1-3u+3u^2$; for B use
$p\le1/2-u+3u^2/2$, $u\le1/2$, and convexity of this quadratic on
$[3/32,1/2]$. Also on S,
\[
\frac{v}{(1-u)^2}\le\frac{2(u-u^2-3/100)}{(1-u)^2}<\frac23.
\]
For two conditional probability vectors, Jensen's inequality for the logarithm gives
\[
H_q((1-u)a,ub)\le h(u)+(1-w)H_q(a)+wH_q(b),\qquad
w=\frac{u^q}{(1-u)^q+u^q}.
\]
One conditional vector in each chart is uniform-two. We claim
\[
H_q(\lambda_S)\le1+h(u)-(1+2u)v,\qquad
H_q(\lambda_B)\le1+h(u)-(3-2u)v.
\tag{CG3}\label{eq:fqe-capped-pairs}
\]
For S, the conditional squared imbalance $z=v/(1-u)^2$ is less than $2/3$.
Its smaller probability is therefore greater than $1/12$; the second bound in
\eqref{eq:fqe-capped-binary} gives a deficit
$z+(3/32)z(1-z)\ge(33/32)z$.
Its weight times $z$ is $v/[(1-u)^2(1+\rho^q)]$ with $\rho=u/(1-u)<1$, and
$\rho^q\le\rho^{3/2}\le\rho(\rho+1/4)$ because $\sqrt\rho\le\rho+1/4$. Hence
$(1-u)^2(1+\rho^q)\le(1-u)^2+u^2+u(1-u)/4$, and the S inequality follows from
\[
 32(1+2u)\Bigl((1-u)^2+u^2+\frac{u(1-u)}4\Bigr)=32+8u(1-7u+14u^2)\le33 .
\]
Indeed $1-7u+14u^2=14(u-1/4)^2+1/8$ is at most $1$ for $u\le1/8$ and at most $11/32$ for
$1/8\le u\le4/15$, so $8u(1-7u+14u^2)\le\max\{1,11/15\}$.

For B the binary cubic bound gives deficit at least $z$ at order $q$.
Here $M=(1-u)^q+u^q\le1-u$, because $q\ge3/2$, $\sqrt u<1/3$, and
$(1-u)^{3/2}\le(1-u)(1-u/2)$.
Thus $u^{q-2}/M\ge u^{-2/5}/(1-u)\ge3-2u$ on $u\le3/32$, as $u^2((1-u)(3-2u))^5\le1$ there:
its logarithmic derivative has the sign of $24u^2-35u+6>0$, and at $u=3/32$ the inequality is
$(3/32)^2(1305/512)^5<1$.
This proves \eqref{eq:fqe-capped-pairs}.

In these coordinates the cubic bound is
$1+4u-4u^2-(1+2u)v$ on S and
$1+4u-4u^2-(3-2u)v$ on B, and the correction is $(11/10)d(1/4-d)$.
By~\eqref{eq:fqe-capped-pairs} the slack is at least $4u(1-u)-h(u)$, which exceeds the correction
by the binary estimate above, since $d\le u-u^2=d_0$. This completes the proof, with strictness.
\end{proof}

\begin{corollary}[The capped cluster region; analytic]
\label{cor:fqe-capped-cluster}
For three four-point spectra with $1\le P\le101/100$, $F\ge11/25$ and every
$p_j\le47/100$, one has $\mathcal E_{q_s}<5/3$.
No quadratic variance bound is required.
\end{corollary}
\begin{proof}
Put $\varepsilon=P-1$. Convexity gives
$U\le(47/100)^2+(7/25+\varepsilon)^2+1/16$.
The sum of the corrections in \eqref{eq:fqe-capped-cubic} is therefore at least
$(11/10)(33/2500+19\varepsilon/100-\varepsilon^2)$.
The cubic envelope gives the uncorrected bound
$3\mathcal E_{q_s}\le5+(1-20\varepsilon-2F)/9$.
Subtracting the correction leaves at most $5-89/75000<5$:
the maximum is at $F=11/25$, $\varepsilon=0$.
\end{proof}


\subsection{An analytic moment bound near the uniform spectrum}\label{sec:sup-nu-moment}
The following argument bounds the near-uniform cuts in the cluster sliver. It applies to all four-point
spectra, without a chart reduction or physical compatibility assumption. Its rectangle form,
Lemma~\ref{lem:fqe-moment-rectangle}, is shared with the bounded-head chart and the capped region.

\begin{lemma}[Near-uniform moment minorant; analytic]
\label{lem:fqe-near-uniform-moment}
Let $\lambda$ be a four-point probability vector. Put
$w_i=\lambda_i-1/4$, $x=\sum w_i^2\le1/25$, and $y=\sum w_i^3$.
For $q=q_s=1587/1000$, write $r=q-1$ and
\[
A=2qr,\quad B=\frac83qr(q-2),\quad
C=\frac83qr(q-2)(q-3),\quad D_5=\frac23(q-4)C.
\]
If $x>0$, set $v=y/x$, $V=A+Bv+Cv^2$, $W=C/4+D_5v$.
The notation $D_5$ distinguishes this Taylor coefficient from the
global third-moment sum $T$.
Then $|v|<3/25$, $V,W>0$, and
\[
2-H_q(\lambda)\ge
\frac{\ln(1+xV+x^2W)}{r\ln2}.
\tag{NU}\label{eq:fqe-near-uniform-moment}
\]
\end{lemma}
\begin{proof}
The sixth derivative of $z^q$ is positive on $z>0$. Taylor's formula
through degree five at $z=1/4$, summed over the four entries, therefore gives
\[
4^{q-1}\sum_i\lambda_i^q
\ge1+Ax+By+C\sum_iw_i^4+D_5xy.
\]
Here the trace-zero Newton identity is
$\sum w_i^5=5xy/6$. The estimate extends to zero entries by continuity.
The elementary square identity
\[
\sum_i\left(w_i^2-\frac x4-\frac yx w_i\right)^2
=\sum_iw_i^4-\frac{x^2}4-\frac{y^2}x\ge0
\]
proves the required moment lower bound, since $C>0$.
Lemma~\ref{lem:fqe-third-moment} gives $|y|\le x^{3/2}/\sqrt3$,
hence $|v|\le1/\sqrt{75}<3/25$. The logarithm is increasing,
and its division by $r\ln2$ gives \eqref{eq:fqe-near-uniform-moment}.
For completeness, the rational bounds at $q_s$
\[
 A\ge\tfrac95,\quad |B|\le\tfrac{21}{20},\quad
 C\ge\tfrac75,\quad -\tfrac{12}5\le D_5<0
\]
give $V\ge\tfrac{837}{500}$ and
$W\ge\tfrac7{20}-\tfrac{12}5\cdot\tfrac3{25}
=\tfrac{31}{500}>0$ for $|v|<3/25$.
This proves the asserted positivity without relying on a later case split.
\end{proof}

The bound~\eqref{eq:fqe-near-uniform-moment} is used three times in the same way: on a rectangle
in $(x,v)$ it is concave enough in $x$ to reduce to the far edge $x=X$, and convex enough in $v$ there
to reduce to one point.
\begin{lemma}[Moment rectangle; analytic]
\label{lem:fqe-moment-rectangle}
Keep $q=q_s$ and the notation of Lemma~\ref{lem:fqe-near-uniform-moment}, put $L=\ln2$,
$N(x,v)=1+xV(v)+x^2W(v)$ and $\rho=21/50=(3/5)(7/10)$, so that $rL<\rho$. For constants $b,c$ and
$\gamma\ge0$ let
\[
\Phi(x,v)=\frac{\ln N(x,v)}{rLx}+bv-c+\gamma x\qquad(x>0).
\]
Let $X>0$ and $v_*\in\mathcal I=[v_{\min},v_{\max}]$, and suppose that on $[0,X]\times\mathcal I$
\[
1\le N\le N_+,\quad V\le V_+,\quad W\ge-W_-,\quad V^2-2W\ge Q,\quad
|B+2Cv+XD_5|\le S,
\tag{MR1}\label{eq:fqe-rect-bounds}
\]
with $W_-\ge0$, and that the constants $m,\kappa,f_0,\eta$ satisfy
\[
Q-2XV_+W_-\ge mN_+^2,\quad m>2\rho\gamma,\quad
2C-XS^2\ge\kappa\rho N_+^2,\quad \kappa>0,
\tag{MR2}\label{eq:fqe-rect-curvature}
\]
\[
\Phi(X,v_*)>f_0,\qquad|\Phi_v(X,v_*)|<\eta.
\tag{MR3}\label{eq:fqe-rect-point}
\]
Then $\Phi$ is strictly decreasing in $x$ and $\Phi>\mu:=f_0-\eta^2/(2\kappa)$ on
$(0,X]\times\mathcal I$. Consequently every four-point probability vector with $0<x\le X$ and
$v=y/x\in\mathcal I$ satisfies
\[
2-H_q(\lambda)-cx+by+\gamma x^2>\mu x.
\tag{MR}\label{eq:fqe-moment-rectangle}
\]
\end{lemma}
\begin{proof}
Fix $v\in\mathcal I$ and put $\ell(x)=\ln N$. Then
\[
\ell''(x)=\frac{2W-V^2-2xW(V+xW)}{N^2}.
\]
Since $N\ge1$, $V+xW\ge0$ for $x>0$. The last term is therefore nonpositive when $W\ge0$; when $W<0$ it is
at most $2XW_-V_+$, as $V+xW\le V$. The numerator is at most $-Q+2XV_+W_-\le-mN_+^2<0$, so
$\ell''\le-m$. Put $\mathcal F(x)=rLx\Phi=\ell(x)-rLx(c-bv-\gamma x)$. Then $\mathcal F(0)=0$ and
$\mathcal F''=\ell''+2rL\gamma\le-m+2\rho\gamma<0$, so $\mathcal F(x)/x$, and with it $\Phi$, is strictly
decreasing in $x$ on $(0,X]$.

At $x=X$, since $1\le N\le N_+$, $C>0$ and $2C-XS^2>0$,
\[
\Phi_{vv}=\frac{2CN-X(B+2Cv+XD_5)^2}{rLN^2}\ge\frac{2C-XS^2}{\rho N_+^2}\ge\kappa .
\]
Completing the tangent square at $v_*$ gives
$\Phi(X,v)\ge\Phi(X,v_*)-\Phi_v(X,v_*)^2/(2\kappa)>\mu$, and so $\Phi(x,v)\ge\Phi(X,v)>\mu$.

The Taylor and projection inequality in the proof of Lemma~\ref{lem:fqe-near-uniform-moment} holds for
every four-point vector; $x\le1/25$ entered only the bound on $|v|$. As $N\ge1>0$, it gives
$2-H_q(\lambda)\ge\ln N/(rL)=x\Phi+cx-by-\gamma x^2$, which is~\eqref{eq:fqe-moment-rectangle}.
\end{proof}

In each application~\eqref{eq:fqe-rect-point} is checked at one point by rational arithmetic. With
$z=XV(v_*)+X^2W(v_*)\in(0,1)$ and an even $k$, $\ln(1+z)\ge\sum_{j=1}^k(-1)^{j+1}z^j/j$; if this sum is
positive, dividing it by $rL_+X$ bounds $\ln N/(rLX)$ below whenever $L<L_+$. The slope
$\Phi_v(X,v_*)=(B+2Cv_*+XD_5)/(rLN(X,v_*))+b$ is monotone in $L$, so it is bounded by substituting the two
rational endpoints of an enclosure of $L$. Table~\ref{tab:fqe-moment-rectangle} lists the three
instances; each application derives~\eqref{eq:fqe-rect-bounds} from its own coefficient bounds.
\begin{table}[htbp]
\centering\small
\begin{tabular}{@{}lccc@{}}
\toprule
 & NU & SHD (large tail) & CG (moment part)\\
\midrule
target & $\Lambda_\theta(c_2)-\Lambda_\theta(\lambda)\ge dx$ & $\Lambda_\theta(\lambda)\le\Lambda_\theta((1/4)^4)$ & \eqref{eq:fqe-capped-cubic}, $p\le21/50$\\
$X$; $\mathcal I$; $v_*$ & $1/25$; $[-3/25,3/25]$; $1/20$ & $3/32$; $[0,9/50]$; $27/200$ & $17/100$; $[-6/25,6/25]$; $0$\\
$b$; $c$; $\gamma$ & $12\beta$; $\alpha+8\beta-\gamma/2$; $\gamma$ (HT) & $9/5$; $35/8$; $159/100$ & $8/3$; $171/40$; $11/10$\\
$N_+$; $V_+$; $W_-$ & $13/12$; $41/20$; $0$ & $6/5$; $19/10$; $3/50$ & $3/2$; $9/4$; $21/100$\\
$Q$; $m$ & $(837/500)^2-4/3$; $6/5$ & $8/3$; $7/4$ & $8/3$; $1$\\
$S$; $\kappa$ & $753/500$; $5$ & $19/10$; $4$ & $213/100$; $2$\\
$f_0$; $\eta$ & $21/2500$; $1/50$ & $1/800$; $1/100$ & $1/250$; $1/20$\\
$k$; $(L_-,L_+)$ & $4$; $(6931/10^4,1733/2500)$ & $8$; $(693147,693148)/10^6$ & $8$; $(693147,693148)/10^6$\\
$\mu$ & $>d=1/125$ & $=1/800-1/80000$ & $=1/250-1/1600$\\
\bottomrule
\end{tabular}
\caption{The three instances of Lemma~\ref{lem:fqe-moment-rectangle}, all at $q_s$ with $\rho=21/50$.
For NU and SHD, $(\alpha,\beta,\gamma)$ is the multiplier of the lemma and the target is the dual
difference $2-H_q-cx+by+\gamma x^2$; for CG the target is the centred form of~\eqref{eq:fqe-capped-cubic}.
The hypotheses are checked by NU-3, SHD-4 and CG-1 through one shared function.}
\label{tab:fqe-moment-rectangle}
\end{table}

\subsection{An analytic uniform-cut dual bound}\label{sec:sup-nu-dual}
\begin{lemma}[Analytic uniform-cut dual bound; NU]
\label{lem:fqe-near-uniform-dual}
At the order $q_s=1587/1000$, with the HT multiplier row, every four-point spectrum
with $p\le29/100$ satisfies
\[
\Lambda_\theta(c_2)-\Lambda_\theta(\lambda)
\ge d\,(p-1/4),\qquad c_2=(1/4)^4.
\tag{ND}\label{eq:fqe-near-uniform-dual}
\]
HT labels the cluster-sliver multiplier row. The order and multipliers are exactly those of the row HT of
Table~\ref{tab:fqe-dual-co}, and $d=1/125$. Equality occurs only at $c_2$.
\end{lemma}
The proof uses two further constants, checked by NU-3: $f_0=21/2500$ and $\eta=1/50$.
\begin{proof}
Put $b=12\beta$, $c=\alpha+8\beta-\gamma/2$ and $x=p-1/4$.
The uniform moment identities give
\[
\Lambda_\theta(c_2)-\Lambda_\theta(\lambda)
=2-H_q(\lambda)-cx+by+\gamma x^2 .
\tag{UD}\label{eq:fqe-uniform-dual-difference}
\]
By Lemma~\ref{lem:fqe-near-uniform-moment}, $x\le X=1/25$ and $|v|<D=3/25$ when $x>0$. We apply
Lemma~\ref{lem:fqe-moment-rectangle} with the NU column of Table~\ref{tab:fqe-moment-rectangle}.
At $q=q_s$ the coefficients satisfy
\[
\frac95\le A\le\frac{19}{10},\quad
-\frac{21}{20}\le B<0,\quad
\frac75\le C\le\frac32,\quad -\frac{12}5\le D_5<0.
\]
These are rational comparisons at the single order $q_s$ (NU-2). Consequently
\[
\frac{837}{500}\le V\le\frac{41}{20},\quad
0<W<\frac23,\quad 1\le N<\frac{13}{12},\quad
|B+2Cv+XD_5|\le\frac{21}{20}+3D+\frac{12}5X=\frac{753}{500},
\]
and $V^2-2W\ge Q=(837/500)^2-4/3$. This is~\eqref{eq:fqe-rect-bounds} with $W_-=0$.
For~\eqref{eq:fqe-rect-curvature}, $Q>(6/5)(13/12)^2$, $2\rho\gamma\le(21/25)(6/5)<6/5$ as
$0\le\gamma\le6/5$, and $C\ge7/5$ gives
$14/5-(1/25)(753/500)^2>5\rho(13/12)^2$.
For~\eqref{eq:fqe-rect-point} take $k=4$ and
\[
\frac{6931}{10000}<L<\frac{1733}{2500}.
\]
The latter follows by summing ten terms of
$L=2\sum_{k\ge0}((2k+1)3^{2k+1})^{-1}$ and bounding the geometric tail.
Direct rational substitution of the multiplier row gives $\Phi(X,v_*)>f_0$ and
$|\Phi_v(X,v_*)|<\eta$ at $v_*=1/20$. Since $\mu=f_0-\eta^2/10>d$,
\eqref{eq:fqe-moment-rectangle} and~\eqref{eq:fqe-uniform-dual-difference} prove
\eqref{eq:fqe-near-uniform-dual}.
At $x=0$ the spectrum is $c_2$; otherwise $d>0$ makes it strict.
\end{proof}

\subsection{An analytic bound on the distinguished cluster cut}\label{sec:sup-sliver}
We now treat the distinguished cut in the cluster sliver. Elementary entropy estimates reduce the
required bound to finitely many polynomial signs, with no interval covering or derivative grid.
Write $q_s=1587/1000$ and let $\theta_s$ be the
HT row of Table~\ref{tab:fqe-dual-co}. Set
\[
\kappa=4\sqrt3/5,\quad X=\tfrac12+\varepsilon-p,\quad
F_c=-1+8\varepsilon+9p-12t.
\]
\begin{lemma}[Analytic cluster sliver; SL]
\label{lem:fqe-cluster-sliver}
At the order $q_s=1587/1000$,
on either ordered chart, suppose
\[
p\ge47/100,\quad X\ge0,\quad 0\le\varepsilon\le1/100,\quad
[F_c-\kappa X,F_c+\kappa X]\cap[11/25,1/2]\ne\varnothing.
\]
For either sign, put
\[
G_{\pm}=\Lambda_{\theta_s}(c_1)-\Lambda_{\theta_s}(\lambda;q_s)
+\ell_{\theta_s}(\varepsilon,F_c\pm\kappa X)
-\ell_{\theta_s}(0,1/2).
\]
Then $G_{\pm}\ge0$, with equality only at
$\lambda=c_1$, $\varepsilon=0$.
\end{lemma}
\begin{proof}
The proof first bounds the two chart coordinates, then replaces the entropy by an explicit
minorant of its deficit. The resulting four polynomial bounds are reduced to one-variable
endpoint signs by removing $\varepsilon$ and using concavity in $v$.

\emph{Coordinates and domains.}
Put $v=\delta^2$. In the two charts use
\[
\begin{aligned}
\lambda_S&=((1-u+\delta)/2,(1-u-\delta)/2,u/2,u/2),\\
\lambda_B&=((1-u)/2,(1-u)/2,(u+\delta)/2,(u-\delta)/2),\\
p&=\tfrac12-u+u^2+v/2,\\
t_S&=\tfrac14-\tfrac34u+\tfrac34u^2+\tfrac34(1-u)v,\qquad
t_B=\tfrac14-\tfrac34u+\tfrac34u^2+\tfrac34uv.
\end{aligned}
\tag{SV}\label{eq:fqe-sliver-coordinates}
\]
In both charts $0\le u<1/16$, $0\le v<1/20$.
Here are explicit domain reductions. Ordering in $S$ gives $u\le1/2$,
$v\le(1-2u)^2$ and $p\le1-3u+3u^2$; hence $u<1/4$.
Since $X\le1/25$ and $\kappa<7/5$, the window condition gives
$(9/2)(1-2u)v\le49/250$.
Successively this gives $v<9/100$, $u<1/10$, $v<11/200$,
$u<1/16$ and $v<1/20$, using $u-u^2\le3/100+v/2$ each time
(verifier-checked, SL-3).
In $B$, ordering gives $u\le1/2$, $v\le u^2$ and
$p\le1/2-u+3u^2/2$; thus already $u<1/30$.

\emph{Entropy estimates.}
Set $(K,c_s)=(7801/2000,5/2)$.
On $0\le u\le1/16$, the binary entropy obeys
\[
H_{q_s}(u,1-u)\le Ku-c_su^{8/5}+\frac{12}5u^2.
\tag{SE}\label{eq:fqe-sliver-binary}
\]
Indeed $(1-u)^q\ge1-qu+q(q-1)u^2/2$. If
$z=qu-q(q-1)u^2/2-u^q$, then $0\le z\le qu$ and
$\ln(1-z)\ge-z-z^2/(2(1-z))$ give the sharper bound
\[
H_q(u,1-u)\le\frac q{(q-1)\ln2}u-\frac{u^q}{(q-1)\ln2}
+\frac{q^2/(2(1-q/16))-q(q-1)/2}{(q-1)\ln2}u^2.
\]
The last coefficient is less than $12/5$. The linear coefficient is
less than $K$, while $16^{8/5-q_s}>51/50$ exceeds $c_s(q_s-1)\ln2$.
These are rational comparisons, integer-power comparisons and the
atanh bounds $693147/10^6<\ln2<693148/10^6$.

Separating the head and tail pairs in the power sum gives
\[
H_{q_s}(\lambda)\le1+Ku-c_su^{8/5}+\frac{12}5u^2
-\begin{cases}
\frac{113}{100}v+\frac32uv-\frac38v^2,&S,\\
3v,&B.
\end{cases}
\tag{SH}\label{eq:fqe-sliver-head-tail}
\]
To see the head estimate, let $M=(1-u)^q+u^q$.
The even Taylor expansion of the head power gives an increase of at least
$(q(q-1)/2)(1-u)^{q-2}v$ over $M$. The factor
\[
\frac{(1-u)^{q-2}}M
=\frac{(1-u)^{-2}}{1+(u/(1-u))^q}\ge1+\frac32u
\]
follows from $(u/(1-u))^q\le7u/25$ and $(1-u)^{-2}\ge1+2u$.
The coefficient $\frac{q(q-1)}2(1-u)^{q-2}/M$ of $v$ in the resulting relative increase is at most
$\frac{q(q-1)}2(1-u)^{-2}\le\frac{q(q-1)}2\bigl(\frac{16}{15}\bigr)^2<27/50$. Applying $\ln(1+z)\ge z-z^2/2$ proves the $S$ estimate,
since $q/(2\ln2)>113/100$ and
$(27/50)^2/(2(q-1)\ln2)<3/8$.
In $B$ the same argument acts on the tail: its factor is
$u^{q-2}/M\ge3$, and the logarithm increment is at most $1/100$
when $v\le u^2$. The resulting coefficient exceeds $3$.
All estimates extend to $u=0$ by continuity.

\emph{Polynomial lower bounds.}
Inserting \eqref{eq:fqe-sliver-head-tail} in $G_{\pm}$
gives the following lower bound, throughout
$u\le1/16$, $v\le1/20$, $\varepsilon\le1/100$:
\[
P(u,v)+\varepsilon(A+Bv),\qquad
P=au+bv+cu^2+duv+fv^2+gu^2v+c_su^{8/5}.
\tag{SP}\label{eq:fqe-sliver-polynomial}
\]
The coefficients are exact rationals. In each row label, the first letter is the chart ($S$ or $B$)
and the second is the sign in $G_\pm$:
\[
\begin{array}{c|rrrrrr|rr}
&a&b&c&d&f&g&A&B\\\hline
% SL-COEFFICIENTS-BEGIN
S-&2/125&7/500&-571/100&-33/10&-281/100&6&185/100&8\\
S+&2/125&7/500&-6&69/50&-103/20&0&147/100&16\\
B-&2/125&19/100&-6&3&-5&0&1/2&0\\
B+&2/125&19/100&-571/100&-3/2&-5/2&0&1/2&0
% SL-COEFFICIENTS-END
\end{array}
\tag{ST}\label{eq:fqe-sliver-coefficients}
\]
For clarity, this table involves polynomial algebra only. Expand the
formula for $G$ using \eqref{eq:fqe-sliver-coordinates}, enclose $\kappa$
between $13856/10000$ and $13857/10000$, and bound negative terms of
higher degree by $u\le1/16$, $v\le1/20$, $\varepsilon\le1/100$.
At $\varepsilon=0$, charge them to $u^2$, $uv$ or $v^2$; in the
coefficient of $\varepsilon$, retain positive terms proportional to $v$.
The difference from every row in \eqref{eq:fqe-sliver-coefficients}
then has nonnegative coefficients (verifier-checked, SL-2). No transcendental sign is left.

\emph{Removing $\varepsilon$ (verifier-checked, SL-3).}
In $S$ the window and $X\ge0$ imply
\[
v\le\min\{2(u-u^2)+2\varepsilon,\ V(u)+21\varepsilon/10\},\qquad
V(u)=\frac{3/50+(7/5)(u-u^2)}{26/5-9u}.
\]
Put $v_0=\min\{v,2(u-u^2),V(u)\}$. Then $v-v_0\le21\varepsilon/10$.
From the two $S$ rows, $P_v\ge-C_0+2fv$, where
$C_0=\max\{0,-b+\max(0,-d)/16\}$. Direct substitution gives
\[
A-\tfrac{21}{10}C_0+\tfrac1{20}\min\{0,B+\tfrac{21}5f\}>9/10.
\]
Thus $G_\pm\ge P(u,v_0)+9\varepsilon/10$. Hence it remains to prove $P\ge0$ with
$\varepsilon=0$, $v\le\min\{2(u-u^2),V(u)\}$ and $u\le1/16$.
In chart $B$ simply discard the positive $\varepsilon/2$ term;
there $u<1/30$ and $v\le u^2$.

\emph{One-variable signs.}
The function $P$ is concave in $v$, so its minimum is at an endpoint.
At $v=0$ the claim follows from $u^{8/5}\ge3u^2$ on $0\le u\le1/16$. In chart $B$,
$b+du+fv\ge3/40$, so this also proves the $v=u^2$ case.

For the $S$ endpoint $v=2(u-u^2)$, that endpoint is the smaller one
only when $u<1/140$. Dividing $P$ by $u>0$ and dropping positive terms gives
\[
P/u\ge a+2b+c_su^{3/5}-Du,\quad
D=2b-c-2\min\{d,(139/140)d\}-4f.
\]
This is a concave function of $u$. It is positive at both $0$ and $1/140$;
for the latter check use $(1/140)^{3/5}>51/1000$.

Finally $v=V(u)$ is the smaller endpoint only for $u>1/150$.
Use the following quadratic minorant of $u^{8/5}$:
for $Z=23/40$, $U=Z^5$, $z=2/5$, $u_0=z^5$, put
\[
Q_{Z,z}(u)=z^8+\frac85z^3(u-u_0)
+\frac{Z^8-z^8-(8/5)z^3(U-u_0)}{(U-u_0)^2}(u-u_0)^2.
\tag{SQ}\label{eq:fqe-sliver-quadratic}
\]
Since the third derivative of $u^{8/5}$ is negative,
Hermite's remainder gives $u^{8/5}\ge Q_{Z,z}(u)$ for $0\le u\le U$.
After replacing $u^{8/5}$ by this minorant, multiply $P(u,V(u))$
by $(26/5-9u)^2$. The resulting polynomial $R$ has degree at most five.
For $L=1/150$, $U=Z^5$, write
\[
R(L+(U-L)x)=\sum_{i=0}^n B_i\binom ni x^i(1-x)^{n-i}.
\]
Every $B_i$ is positive. Explicit integer lower bounds on $10^6 B_i$ are
\[
\begin{array}{c|r|rrrrrr}
&n&\multicolumn{6}{c}{\text{coefficient lower bounds}}\\\hline
% SL-BERNSTEIN-BEGIN
S-&5&793&27467&39030&38325&27813&9606\\
S+&4&238&51693&86240&106860&116069&
% SL-BERNSTEIN-END
\end{array}
\tag{SB}\label{eq:fqe-sliver-bernstein}
\]
These are two polynomial expansions on single intervals, with no
subdivision. They prove the last endpoint signs. Together with the
strict $v=0$ and $\varepsilon$ bounds, this proves the asserted equality case.
\end{proof}


\subsection*{Further results}
\noindent The proof of Corollary~\ref{cor:fqe-renyi-gw} in the paper uses none of the results of
Sections~\ref{sec:fqe-renyi-above-one}--\ref{sec:fqe-renyi-filter-limit}. They give an independent near-one defect estimate
with its quantitative form, the sharp purity-variance profile on the critical locus with an analytic proof of (QV) on critical
states (used for the equality case of (QK) at $P=1$), and a shorter proof of the envelope for states with a rank-two cut.
\subsection{An independent near-one defect estimate}\label{sec:fqe-renyi-above-one}

Use $H_q$, $h_q$ and $\mathcal E_q$ from Section~\ref{sec:fqe-spectral}, and set
$m_q=2^{-q}+3\,6^{-q}$. Entropies are in bits throughout this section; the proof below uses
$H_q^{\mathrm{nat}}=H_q\ln2$ when working in nats.
The Tsallis result alone does not give an upper R\'enyi bound when $q>1$.
The anchor-order argument proves the full HS--cluster envelope of Corollary~\ref{cor:fqe-renyi-gw} for all $q>1$.
Here we give an independent near-one proof using the nonnegative Hermite interpolation error to control
the logarithm's curvature; it also yields the quantitative defect estimate~\eqref{eq:fqe-renyi-defect}.

\begin{theorem}[R\'enyi interval above one]
\label{thm:fqe-renyi-above-one}
For every pure four-qubit state and every $1<q\le5/4$,
\[
\boxed{\quad\mathcal E_q(\psi)\le\frac{\log_2 m_q}{1-q}.\quad}
\]
Equality holds exactly on the two $M_4$, $\overline{M_4}$ LU orbits.
\end{theorem}


The argument rests on a scalar slack estimate. Set $p=1/6$, $s=1/2$, and
\[
f_q(x)=\frac{x^q-x}{1-q},\qquad f_1(x)=-x\ln x,\qquad f_q(0)=0.
\]
Let $\ell_q$ be the secant of $f_q$ through $p,s$, and $\widehat h_q$ its cubic
Hermite interpolant there, as in Corollary~\ref{cor:fqe-tsallis}.

\begin{lemma}[Scalar slack control]
\label{lem:fqe-renyi-slack}
For $1\le q\le5/4$ and $0\le x\le1$,
\[
 \bigl(f_q(x)-\ell_q(x)\bigr)^2
 \le\frac{141}{100}\bigl(\widehat h_q(x)-f_q(x)\bigr).
 \tag{RA}\label{eq:fqe-renyi-scalar}
\]
Consequently, for any four-component probability vector $\lambda$, put
$S_q=\sum_i f_q(\lambda_i)$, $S_q^*=f_q(s)+3f_q(p)$,
and $D_q=\sum_i(\widehat h_q-f_q)(\lambda_i)$. Then
\[
 (S_q-S_q^*)^2\le\frac{141}{25}D_q.
 \tag{RB}\label{eq:fqe-renyi-sum}
\]
\end{lemma}
The proof (Section~\ref{sec:sup-renyi-slack}; check RA-1) puts
$G=\widehat h_q-f_q-\frac{100}{141}(f_q-\ell_q)^2$, which vanishes with its derivative at $p$ and $s$, and shows
$G^{(4)}>0$ on $(0,1]$ from an elementary bound $\Phi_q<141/100$ on the factor its fourth derivative contains;
the Hermite remainder then gives $G\ge0$, and Cauchy--Schwarz gives the summed bound.

\begin{proof}[Proof of Theorem~\ref{thm:fqe-renyi-above-one}]
Put $r=q-1$, $L=4^{-r}$, $M=\sum_i\lambda_i^q=1-rS_q$,
and $m=m_q=1-rS_q^*$. Convexity of $x^q$ gives $M\ge L$ and $m\ge L$.
For $u>0$, integration of $(t-1)/t$ gives
\[
u-1-\ln u\le\frac{(u-1)^2}{2\min(u,1)}.
\]
Apply this with $u=M/m$; since $\min(u,1)\ge L/m$, it follows that
\[
H_q^{\mathrm{nat}}(\lambda)-H_q^{\mathrm{nat}}(s,p,p,p)
\le\frac{S_q-S_q^*}{m}+
\frac{r(S_q-S_q^*)^2}{2mL}.
\]
By \eqref{eq:fqe-renyi-sum} and $S_q=\sum_i \widehat h_q(\lambda_i)-D_q$, this is at most
\[
\frac{\sum_i \widehat h_q(\lambda_i)-S_q^*}{m}
-\frac{1-141r/(50L)}{m}D_q.
\tag{RD}\label{eq:fqe-renyi-cut}
\]
For $r\le1/4$, the last coefficient is positive, because
\[
\frac{141r}{50L}\le\frac{141}{200}\sqrt2<1
\qquad\left(\left(\tfrac{141}{200}\sqrt2\right)^2=\tfrac{19881}{20000}\right).
\]
Write $a_q,b_q$ for the cubic and quadratic coefficients of $\widehat h_q$ and
$\alpha_q=-b_q-17a_q/12$. Corollary~\ref{cor:fqe-tsallis} proves
$a_q,\alpha_q>0$. Sum \eqref{eq:fqe-renyi-cut} over the three cuts and
use $T=(17P-12-F)/12$. With $\mathcal D_q=\sum_jD_q(\rho_{1j})$ and
$\kappa=1-141\sqrt2/200>0$, the result is the defect estimate
\[
3m_q\ln2\bigl(\mathcal E_q(M_4)-\mathcal E_q(\psi)\bigr)
\ge\alpha_q(P-1)+\frac{a_q}{12}F+\kappa\mathcal D_q\ge0.
\tag{RE}\label{eq:fqe-renyi-defect}
\]
Equality forces $P=1,F=0$, so Proposition~\ref{prop:fqe-moment-zeros} gives the two HS
orbits; both attain the bound.
\end{proof}

\subsection{Proof of Lemma~\ref{lem:fqe-renyi-slack}}\label{sec:sup-renyi-slack}
Notation is that of Section~\ref{sec:fqe-renyi-above-one}: $p=1/6$, $s=1/2$,
$f_q(x)=(x^q-x)/(1-q)$, $\ell_q$ the secant of $f_q$ through $p,s$, and $\widehat h_q$ its cubic Hermite
interpolant there.


The proof rests on an elementary bound for the function $\Phi_q$ that appears in the fourth derivative.

\begin{lemma}[An elementary fourth-derivative bound; proved, RA-1]
\label{lem:fqe-slack-fourth-derivative}
Write $\ell_q(x)=\mu_qx+\nu_q$ and, for $1<q\le5/4$, put
\[
 A(q)=\frac{2(q+1)}{3-q}\left(\frac1{q-1}-\mu_q\right),\qquad
 D(q)=\frac{4(2q-1)(3-2q)}{(q-1)(3-q)(2-q)},
\]
and $\Phi_q(x)=2\nu_q+A(q)x-D(q)x^q$.
Define $\Phi_1$ by continuity. Then $\Phi_q(x)<141/100$ for
$1\le q\le5/4$ and $0\le x\le1$.
\end{lemma}
\begin{proof}
Set $r=q-1$, $p=1/6$, $s=1/2$, and $\delta=r(2-r)(1-r)$.
The function $\Phi_q$ is concave in $x$. Its derivative at $p$ is
$p^rN_p/\delta$, where
\[
 N_p=(2+r)(1-r)(3\cdot3^r-1)-4(1+r)(1-4r^2).
\]
Since $\ln3>1$, $3^r\ge1+r$, and substitution gives
$N_p\ge11r^2+13r^3\ge0$. Thus its maximum on $[0,1]$ is attained
in $[p,1]$.

We show that $\Phi_q(x)$ increases with $q$ on this latter interval.
Three endpoint derivatives suffice. Define
\[
 \begin{split}
 N_0&=6(1-r)3^r-6+5r+7r^2,\\
 N_1&=(2+r)(1-r)(3\cdot2^{-r}-6^{-r})-4(1+r)(1-4r^2),\\
 J_0&=(N_0'-(\ln6)N_0)\delta-N_0\delta',\\
 J_p&=(N_p'-(\ln6)N_p)\delta-N_p\delta',\qquad
 J_1=N_1'\delta-N_1\delta'.
 \end{split}
\]
Primes here mean $r$ derivatives. Direct substitution into $\Phi_q$ gives
\[
 \partial_q\Phi_q(p)=\frac{p^rJ_0}{3\delta^2},\qquad
 \partial_q\partial_x\Phi_q(p)=\frac{p^rJ_p}{\delta^2},\qquad
 \partial_q\partial_x\Phi_q(1)=\frac{J_1}{\delta^2}.
\]
All three $J$ are positive for $0<r\le1/4$. Here are coarse bounds
that permit this assertion to be checked with small rational numbers.
Write $J=\sum_{k=2}^6c_kr^k+R$; the constant and linear coefficients
are zero. The five coefficients have the following lower bounds,
and $|R|\le Mr^7$:
\[
\begin{array}{c|rrrrr|r}
 &c_2&c_3&c_4&c_5&c_6&M\\\hline
% RA-TAYLOR-BEGIN
J_0&4&-6&15&-14&0&3\\
J_p&27&-2&-33&52&-29&3\\
J_1&4&86&-75&2&0&14\\
% RA-TAYLOR-END
\end{array}
\]
For completeness, these bounds require only
$69/100<\ln2<7/10$ and $109/100<\ln3<11/10$.
Both follow from the first five terms of
$\ln u=2\sum_{j\ge0}z^{2j+1}/(2j+1)$, $z=(u-1)/(u+1)$,
with tail at most $2z^{11}/[11(1-z^2)]$.
Expand the displayed $N$ and $J$ through degree six and substitute these
two logarithm intervals to get the five columns. To check the remainder,
each $J$ is a polynomial of degree at most six plus terms
$Q(r)e^{\alpha r}$ with $\deg Q\le5$.
For $Q(r)=\sum_jQ_jr^j$, its tail after degree six has absolute value at most
\[
 r^7 e^{|\alpha|/4}\sum_j
       \frac{|Q_j|\,|\alpha|^{7-j}}{(7-j)!}.
\]
Use $3^{1/4}<4/3$, $6^{1/4}<8/5$ and $2^{1/4}<6/5$ for the
three possible exponential factors. The resulting sums are respectively
less than $3,3,14$, giving the last column. In particular the table yields
$J_0>2r^2$, $J_p>24r^2$ and $J_1>3r^2$ on $0<r\le1/4$;
for the last row use $86r-75r^2\ge0$.

Now let $W(x)=\partial_q\partial_x\Phi_q(x)$ and
$k(q)=q(q-1)D(q)>0$. Differentiating $\partial_x^2\Phi_q=-k(q)x^{q-2}$ gives
\[
 W'(x)=-k(q)x^{q-2}\bigl(\partial_q\ln k(q)+\ln x\bigr).
\]
The bracket increases with $x$, so $W$ has no interior minimum.
Its endpoint values at $p,1$ are positive by the table. Thus $W\ge0$,
and $\partial_q\Phi_q(x)\ge\partial_q\Phi_q(p)>0$ on $[p,1]$.
It remains only to bound $\Phi_{5/4}$.

Put $a=6^{-1/4}$, $b=2^{-1/4}$, and $t=x^{1/4}$. Then
\[
 \Phi_{5/4}(x)=2(b-a)+\frac{36}{7}(3b-a)t^4-\frac{64}{7}t^5
 \le2(b-a)+\frac{36\cdot9^4}{35\cdot20^4}(3b-a)^5,
\]
by maximising the last two terms over $t\ge0$; their maximum is at
$t=9(3b-a)/20$. Integer fourth powers give $a>319/500$ and $b<841/1000$.
Consequently
\[
 \Phi_{5/4}(x)
 <\frac{203}{500}+\frac{36\cdot9^4}{35\cdot20^4}
                         \left(\frac{377}{200}\right)^5
 <\frac{141}{100}.
\]
Concavity in $x$ and monotonicity in $q$ give the assertion everywhere;
the endpoint $q=1$ follows by continuity.
\end{proof}

\begin{proof}[Proof of Lemma~\ref{lem:fqe-renyi-slack}]
Let $g=f_q-\ell_q$ and
$G=\widehat h_q-f_q-(100/141)g^2$. Both $G$ and $G'$ vanish at $p,s$.
Since $f_q^{(4)}=-q(2-q)(3-q)x^{q-4}<0$, differentiation gives
\[
 G^{(4)}(x)=-f_q^{(4)}(x)\left(1-\frac{100}{141}\Phi_q(x)\right)>0
 \qquad(x>0).
\]
The cubic Hermite remainder, with repeated nodes $p,p,s,s$, gives
$G(x)=G^{(4)}(\xi)(x-p)^2(x-s)^2/24\ge0$ for $x>0$;
continuity gives the result at zero. No derivative at zero is needed.
Finally $\sum_i\ell_q(\lambda_i)=S_q^*$, since $\ell_q$ is affine and
$\sum_i\lambda_i=s+3p=1$. Cauchy--Schwarz gives the summed bound.
\end{proof}

\subsection{Critical-locus variance refinements}\label{sec:fqe-renyi-compatibility}
The universal bound (QK) suffices for the anchor argument. We now give its sharper
critical-locus refinements: the fixed-variance profile on $P=1$, used to classify equality
in (QK), and the analytic proof of the sharp bound~\eqref{eq:fqe-quadratic-variance} on every critical state.
Put $p_j=\operatorname{tr}\rho_{1j}^2$ and $U=\sum_jp_j^2$.

\begin{proposition}[Sharp purity-variance tradeoff at minimal purity]
\label{prop:fqe-critical-variance}
Suppose a pure four-qubit state has $P=1$. Set
\[
V=\sum_{j=2}^4(p_j-1/3)^2,\qquad s=(24V)^{1/4}.
\]
Then $0\le s\le1$ and the sharp bound at each fixed $V$ is
\[
\boxed{\quad F\ge\frac32s^2-s^3.\quad}
\tag{RF}\label{eq:fqe-critical-variance}
\]
For $a=(1+2s)/3$, $b=(1-s)/3$, equality holds exactly when, up to ordering of
cuts and eigenvalues, the three spectra are
\[
\begin{split}
&\left(\frac12,\frac a2,\frac b2,\frac b2\right),\\
&\left(\frac{1+b+2\sqrt{ab}}4,\frac{1+b-2\sqrt{ab}}4,
\frac{1-b}4,\frac{1-b}4\right)\quad\text{on each of the other two cuts}.
\end{split}
\tag{RF-eq}\label{eq:fqe-variance-spectra}
\]
These spectra are attained for every $s\in[0,1]$.
In particular, $T+U\le3/4$, equivalently $F\ge12V$; equality in this weaker
bound occurs only at the HS spectra ($s=0$) or cluster spectra ($s=1$).
\end{proposition}
\noindent[Proved under the hypothesis $P=1$.]
\begin{proof}
By Lemma~\ref{lem:fqe-purity}, the state is critical with $\tau=0$. As in the proof of
Proposition~\ref{prop:fqe-critical-envelope}, its Bloch correlation matrix is then $\Theta=\mathbf m\mathbf n^T$ with unit
$\mathbf m,\mathbf n\in\mathbb R^3$ (Lemma~\ref{lem:fqe-critical-correlation}(2) at $\kappa=0$), and the cut $c\in\mathbb Z/3$ has the
four eigenvalues $(1+\sigma\cdot g^{(c)})/4$, $g^{(c)}_a=m_{a+c}n_a$, as $\sigma$ runs over the sign triples with
$\sigma_1\sigma_2\sigma_3=-1$ (Lemma~\ref{lem:fqe-p1-cuts}). Put $A_a=m_a^2$, $B_a=n_a^2$ and $k=-m_1m_2m_3n_1n_2n_3$.
The three vectors $g^{(c)}$ have squared norms
\[
u_c=\sum_aA_{a+c}B_a,\qquad \sum_cu_c=1,
\]
and the same coordinate product $-k$. The sign sums $\sum_\sigma\sigma\cdot g=0$, $\sum_\sigma(\sigma\cdot g)^2=4|g|^2$ and
$\sum_\sigma(\sigma\cdot g)^3=-24g_1g_2g_3$ give, for the corresponding cut,
\[
p_j=\frac{1+u_c}{4},\qquad
\operatorname{tr}\rho_{1j}^3=\frac{1+3u_c+6k}{16}.
\]
Set $X=1-3e_2(A)$ and $Y=1-3e_2(B)$, where $e_2(A)=\sum_{i<j}A_iA_j$.
Then $0\le X,Y\le1$, and expanding the three cyclic sums gives
\[
\sum_c u_c^2=\frac{1+2XY}{3}.
\]
We sharpen the elementary Schur estimate by keeping $X$ and $Y$.
For any nonnegative triple $A$ of sum one and $X=1-3e_2(A)$,
\[
27e_3(A)\le h(X),\qquad
h(X)=1-3X+2X^{3/2}=(1-\sqrt X)^2(1+2\sqrt X).
\tag{RP}\label{eq:fqe-triple-product}
\]
This is the cubic discriminant identity
\[
 4X^3-\bigl(27e_3(A)-1+3X\bigr)^2
   =27\prod_{i<j}(A_i-A_j)^2\ge0.
\]
Taking the upper square-root bound proves~\eqref{eq:fqe-triple-product}.
Equality requires two equal coordinates, and the upper sign selects
$((1+2r)/3,(1-r)/3,(1-r)/3)$, $r=\sqrt X$, up to permutation.
This includes the uniform triple and coordinate-vector endpoints.

The transfer from $X,Y$ to their geometric mean is algebraic as well.
Set $g(t)=(1-t)^2(1+2t)$, $\xi=X^{1/4}$, $\zeta=Y^{1/4}$,
$q=\xi\zeta$, $d=(\xi-\zeta)^2$, and
$w=(1-\xi^2)(1-\zeta^2)$. Direct factorisation gives
\begin{equation}\label{eq:fqe-product-transfer-remainder}
 \begin{split}
 h(\sqrt{XY})^2-h(X)h(Y)
 &=g(\xi\zeta)^2-g(\xi^2)g(\zeta^2)\\
 &=d\bigl\{w[6q(q+2)+2d]+d(1+2q)^2\bigr\}\ge0.
 \end{split}
\end{equation}
Every factor is nonnegative for $0\le X,Y\le1$, and the difference
is strictly positive when $X\ne Y$. Consequently
\[
27k\le27\sqrt{e_3(A)e_3(B)}
\le\sqrt{h(X)h(Y)}\le h(\sqrt{XY}).
\]
The moment formulas above give
\[
V=\frac{XY}{24},\qquad F=\frac12-\frac{27}{2}k.
\]
Thus $s=(XY)^{1/4}$ and $F\ge(1-h(s^2))/2=3s^2/2-s^3$.
For $0<s<1$, equality requires $X=Y=s^2$, $k>0$, and both $A,B$ to be
permutations of $(a,b,b)$. The shift $c$ that aligns their distinguished
coordinates gives $|g^{(c)}_a|=(a,b,b)$ up to order; the other two shifts give
$(\sqrt{ab},\sqrt{ab},b)$. As every $g^{(c)}$ has negative coordinate product,
all sign choices give the same spectral multisets, and the eigenvalue formula
proves \eqref{eq:fqe-variance-spectra}.
At $s=0$, equality forces $k=1/27$ and hence both triples uniform; at $s=1$
both are coordinate vectors. These give the stated endpoint spectra.

For attainability, take $\mathbf m=-\mathbf n=(\sqrt a,\sqrt b,\sqrt b)$, so $A=B=(a,b,b)$ and $k=ab^2=e_3(A)$.
The product auxiliary vector with Bloch vectors $\mathbf m,\mathbf n$ has $\Theta=\mathbf m\mathbf n^T$ and $\tau=0$, and it is
$\sum_kv_km_k$ for a critical $G(v)$, as at the end of the proof of Proposition~\ref{prop:fqe-critical-envelope}.
Finally,
\[
\frac32s^2-s^3-12V=\frac{s^2}{2}(1-s)(3+s)\ge0,
\]
with equality exactly at $s=0,1$. This proves the weaker bound and its
exhaustive equality statement as a corollary.
\end{proof}

\subsection{A joint-cut proof on the critical locus}
\label{sec:fqe-critical-quadratic-variance}

\begin{proposition}[Sharp quadratic variance on critical states]
\label{prop:fqe-critical-quadratic-variance}
Every critical unit four-qubit state satisfies
\[
 F^2+3P(P-1)\ge6V_P.
 \tag{CQ}\label{eq:fqe-critical-quadratic-variance}
\]
Equality holds at HS spectra and along the full cluster-to-Bell family
\[
 \bigl((1+\kappa)/2,(1-\kappa)/2,0,0\bigr),\qquad
 (1/4)^4,\qquad(1/4)^4,\qquad 0\le\kappa\le1,
\]
up to ordering of cuts and eigenvalues, and only there.
In particular, the defect in \eqref{eq:fqe-quadratic-variance-proved} is at least
$3\tau(1-\tau)/4$ on the critical locus. The proof uses no variance certificate. At $P=1$ the sharper
profile~\eqref{eq:fqe-critical-variance} of Proposition~\ref{prop:fqe-critical-variance} holds.
\end{proposition}
\begin{proof}
Use the auxiliary correlation matrix of Proposition~\ref{prop:fqe-critical-envelope}:
\[
 \Theta=U\operatorname{diag}(\kappa,-\kappa,1)V^T,\qquad U,V\in SO(3),
 \qquad \tau=\kappa^2,\qquad P=1+\kappa^2/2.
\]
Write $r=Ue_3$, $s=Ve_3$, $A_i=r_i^2$, $B_i=s_i^2$, and
$X=1-3e_2(A)$, $Y=1-3e_2(B)$. Thus $0\le X,Y\le1$.
We retain these two directional parameters in both the permanent and purity estimates.

\emph{The permanent bound with directions retained.}
Put $z=\sqrt{XY}$ and
\[
 b_\kappa(x)=\kappa^2+\frac23\kappa(1-\kappa)^2(1-x)
             +\frac29(1-\kappa)^3(1-3x+2x^{3/2}).
\]
The polar-decomposition Cauchy--Schwarz bound from
Proposition~\ref{prop:fqe-critical-envelope}, now using
\eqref{eq:fqe-triple-product} instead of discarding $X,Y$, gives
\[
 |\operatorname{perm}\Theta|
 \le\sqrt{b_\kappa(X)b_\kappa(Y)}\le b_\kappa(z).
\]
For completeness, the last inequality is geometric concavity.
Write $b_\kappa(x)=a-bx+cx^{3/2}$, where
\[
 a=(4\kappa^3+3\kappa^2+2)/9,\quad
 b=2(1-\kappa)^2/3,\quad c=4(1-\kappa)^3/9.
\]
The sign of $d^2\log b_\kappa(e^v)/dv^2$ is the negative of
$4ab-9ac t+bc t^3$, $t=\sqrt x$.
This polynomial decreases on $[0,1]$, since $b\le3a$, and its value at one is
$4\kappa^2(1-\kappa)^2(4\kappa^2+\kappa+1)/9\ge0$.
The cases with zero arguments follow by continuity.
Consequently the exact formula for $F$ gives
\[
 F\ge\kappa^2(1-\kappa)+\frac32(1-\kappa)^2z
                  -(1-\kappa)^3z^{3/2}
 \ge f_\kappa(z),
\]
where, using $\sqrt z\le(1+z)/2$,
\[
 f_\kappa(z)=(1-\kappa)\kappa^2
       +\frac12(1-\kappa)^2(1+2\kappa)z
       +\frac12(1-\kappa)^3z(1-z)\ge0.
\]

\emph{The three purities as one Fourier coefficient.}
Let $\omega=e^{2\pi i/3}$, $D=\operatorname{diag}(1,\omega,\omega^2)$, and
\[
 u_c=\sum_{j=0}^2\Theta_{j+c,j}^2,\qquad
 Z=\sum_c\omega^c u_c
   =\operatorname{tr}(D\Theta\overline D\Theta^T).
\]
Indices are cyclic. The cut formula $p_c=(1+u_c)/4$ implies
$6V_P=|Z|^2/4$.
The symmetric unitary matrices $M=U^TDU$ and $N=V^T\overline D V$
have $|M_{33}|^2=X$, $|N_{33}|^2=Y$.
Their last-row off-diagonal squared norms are $1-X,1-Y$;
their leading $2\times2$ blocks have squared Frobenius norms $1+X,1+Y$.
Expanding $Z$ in these blocks and applying Cauchy--Schwarz gives
\[
 |Z|\le z+2\kappa\sqrt{(1-X)(1-Y)}
                +\kappa^2\sqrt{(1+X)(1+Y)}
       \le B_\kappa(z):=(1-\kappa)^2z+2\kappa+\kappa^2.
\]
To see the last step, fix $z$ and put $S=X+Y\ge2z$.
The last two terms become
$2\kappa\sqrt{1-S+z^2}+\kappa^2\sqrt{1+S+z^2}$,
which decreases with $S$ for $0\le\kappa\le1$; evaluate at $S=2z$.

\emph{An elementary scalar comparison.}
For $0\le\kappa,z\le1$,
\[
 4f_\kappa(z)^2+9\kappa^2-B_\kappa(z)^2
       -3\kappa^2(1-\kappa^2)
   =(1-\kappa)^2(1-z)R_\kappa(z),
\]
where, putting $d=1-\kappa$,
\[
 R_\kappa(z)=-d^4z^3+3d^3(1+\kappa)z^2
                    -4\kappa d^2z+2\kappa^2(1+2\kappa^2).
\]
Since $z^3\le z^2$,
\[
 R_\kappa(z)\ge
 2d^3(1+2\kappa)z^2-4\kappa d^2z+2\kappa^2(1+2\kappa^2)\ge0.
\]
The final quadratic is nonnegative: its discriminant is
$-16\kappa^3d^3(3+2\kappa+4\kappa^2)\le0$ and its leading and constant
coefficients are nonnegative (the endpoints follow directly).
Combining the estimates proves
\[
 4F^2+9\tau-24V_P\ge3\tau(1-\tau),
\]
which is exactly~\eqref{eq:fqe-critical-quadratic-variance}.

For $0<\kappa<1$ the scalar comparison is strict unless $z=1$.
Then $X=Y=1$, and equality in the Fourier triangle bound forces all nonzero
entries of $\Theta$ onto one cyclic diagonal. Their magnitudes are
$1,\kappa,\kappa$, with negative product since $\det\Theta=-\kappa^2$.
The cut eigenvalue formula gives precisely the displayed family.
At $\kappa=0$, equality requires $z=0$ or $z=1$; the latter is its cluster
endpoint, while the former requires $F=0$ and hence HS by the critical
envelope equality case. At $\kappa=1$, equality requires $F=0$ and
$|Z|=3$, forcing the Bell endpoint.
Conversely the family spectra are attained by the correlation matrices
$\operatorname{diag}(\kappa,-\kappa,1)$, and direct substitution gives equality.
\end{proof}



\subsection{Independent arguments on restricted domains}\label{sec:fqe-renyi-restricted}
The proof of Corollary~\ref{cor:fqe-renyi-gw} in the paper needs no eigenvalue cap and no rank condition. The following result gives a shorter,
independent proof of the envelope for states with a rank-two cut, at every order (for $q\ge2$ see
Section~\ref{sec:fqe-renyi-high-order}).

\begin{proposition}[Sharp fixed-spectrum rank-two cut envelope]
\label{prop:fqe-renyi-rank-two}
Suppose one of the three balanced cuts of a pure four-qubit state has
Schmidt rank at most two and nonzero eigenvalues $a,1-a$,
where $1/2\le a\le1$. Then for every $q>0$,
\[
\boxed{\qquad
\mathcal E_q(\psi)\le
\frac43+\frac{1}{3}H_q(a,1-a)
\le\frac53.\qquad}
\tag{Rk2}\label{eq:fqe-renyi-rank-two}
\]
For each fixed $a$, the first bound is sharp; equality holds precisely
when the other two balanced cut spectra are both $(1/4)^4$.
The attained cluster-to-Bell family \eqref{eq:fqe-quadratic-family}
realises equality for every $a\in[1/2,1]$ and every order $q>0$.
In particular the global HS--cluster R\'enyi envelope holds on
this entire rank-two boundary without any purity or other cap hypothesis.
\end{proposition}
\begin{proof}
A density matrix on a four-dimensional cut has R\'enyi entropy
at most $\log_2 4=2$, with equality exactly for the uniform-four
spectrum. The stipulated cut has entropy $H_q(a,1-a)$, so bounding
each of the other two cut entropies by $2$ proves the first
inequality and its equality classification. A binary probability
vector has entropy at most $1$, proving the second inequality.
The physical states in \eqref{eq:fqe-quadratic-family} have the
stipulated binary spectrum and both crossed cuts uniform-four
for all $a=(1+r)/2\in[1/2,1]$; they therefore attain the bound
at every $q>0$. Since the HS entropy $h_q$ is at least $5/3$
up to $q_*$, and the cluster ceiling is $5/3$ thereafter,
the final assertion follows.
\end{proof}


\subsection{A cap-free purity-controlled R\'enyi bound}\label{sec:fqe-renyi-purity-bound}
The defect estimate of Theorem~\ref{thm:fqe-renyi-above-one} also gives a ceiling at every order $q\ge5/4$ that
improves linearly with the purity sum. It is not used in the proof of Corollary~\ref{cor:fqe-renyi-gw}.

\begin{corollary}[Cap-free purity-controlled R\'enyi bound; exact-rational check, RA-5]
\label{cor:fqe-renyi-purity}
For every unit pure four-qubit state and every $q\ge5/4$,
\[
\boxed{\qquad \mathcal E_q(\psi)\le h_{5/4}-\frac{13}{20}(P-1),
\qquad h_{5/4}=\frac{\log_2m_{5/4}}{1-5/4}.\qquad}
\tag{RH}\label{eq:fqe-renyi-high-purity}
\]
In particular, $P\ge10/9$ implies $\mathcal E_q(\psi)<5/3$
for $q\ge5/4$. The bound \eqref{eq:fqe-renyi-high-purity} holds for every triple of four-point spectra with
$P\ge1$ and $F=17P-12T-12\ge0$, since the proof uses only these and the per-cut estimate
\eqref{eq:fqe-renyi-cut}.
\end{corollary}
\begin{proof}
At $q_0=5/4$, the defect estimate \eqref{eq:fqe-renyi-defect}
and $F,\mathcal D_{q_0}\ge0$ give
\[
\mathcal E_{q_0}(\psi)\le h_{q_0}-c(P-1),
\qquad c=\frac{\alpha_{q_0}}{3m_{q_0}\ln2}.
\]
The exact-rational enclosures RA-5 establish
\[
20\alpha_{q_0}-39m_{q_0}\ln2>0,\qquad
4\ln m_{q_0}+\frac{313}{180}\ln2>0.
\]
Equivalently $c>13/20$ and $h_{q_0}<313/180$.
R\'enyi entropy is nonincreasing in its order, so
$\mathcal E_q(\psi)\le\mathcal E_{q_0}(\psi)$ for $q\ge q_0$.
Since $P\ge1$, the first assertion follows. If $P\ge10/9$,
\[
\mathcal E_q(\psi)<\frac{313}{180}
-\frac{13}{20}\left(\frac{10}{9}-1\right)=\frac53,
\]
as asserted.
\end{proof}

\subsection{The unrestricted joint purity bound}\label{sec:fqe-joint-purity}

\begin{theorem}[Joint purity compatibility; consequence of QK]
\label{thm:fqe-joint-purity}
Every pure four-qubit state satisfies
\[
\boxed{\qquad T+\sum_{j=2}^4p_j^2\le3P-\frac94.\qquad}
\tag{JV}\label{eq:fqe-joint-purity}
\]
Equality holds exactly for HS cut spectra, cluster cut spectra, or products
of two Bell pairs, up to ordering of the cuts and eigenvalues.
In the HS case the states belong to the two HS LU orbits.
\end{theorem}
\begin{proof}
This is an elementary consequence of the stronger quadratic variance bound
\eqref{eq:fqe-quadratic-variance-proved}, together with $F\ge0$ and
$p_j\in[1/4,1]$. Set
\[
U=\sum_jp_j^2,\qquad G=3P-\frac94-T-U,\qquad
Q=F^2+\frac92(P-1)-6\bigl(U-P^2/3\bigr)\ge0.
\]
Since $T=(17P-12-F)/12$, direct algebra gives the useful identity
\[
12G=F(1-2F)+2(P-1)(3-2P)+2Q.
\tag{JI}\label{eq:fqe-joint-implication}
\]
It proves $G\ge0$ whenever $1\le P\le3/2$ and $0\le F\le1/2$.

For the remaining cases, use the purity box \eqref{eq:fqe-purity-box} of Section~\ref{sec:sup-qk}, the
maximum of the convex function $U$ at fixed $P$ over $[1/4,1]^3$.
Using $12G=19P-15+F-12U$, the first piece, for
$1\le P\le3/2$ and $F\ge1/2$, yields
\[
12G\ge F-\frac12+(P-1)(19-12P)\ge0.
\]
For $3/2\le P\le9/4$ the second piece yields
\[
12G\ge F+\frac12(2P-3)(31-12P)\ge0,
\]
and for $9/4\le P\le3$ the last piece yields
\[
12G\ge F-12P^2+67P-87\ge F+3>0.
\]
The last inequality follows from concavity of the quadratic and its
values $3,6$ at the two endpoints.

Equality in the first case forces $Q=0$, so the equality classification of
Corollary~\ref{thm:fqe-quadratic-variance} gives HS, cluster or Bell-pair
spectra. In the $F\ge1/2$ case equality requires
$P=1$, $F=1/2$, $U=3/8$, again giving $Q=0$ and the cluster
spectra. In the second box range equality requires
$P=3/2$, $F=0$ and $U=9/8$; its purity triple is
$(1,1/4,1/4)$, and the state is a product of two Bell pairs.
The final range is strict. All three listed types attain equality.
\end{proof}


\subsection{What local filtering can and cannot reduce}
\label{sec:fqe-renyi-filter-limit}

\begin{lemma}[One-site filtering majorises a locally maximally mixed pair]
\label{lem:fqe-single-filter}
Let $\rho$ be a two-qubit density matrix with both one-qubit
marginals equal to $I/2$. Let $K$ be any nonzero filter on
one of its qubits, and normalise
$\rho_K=(K\otimes I)\rho(K^\dagger\otimes I)$ by its trace
(with the factors interchanged if the filter acts on the second qubit).
Then
\[
\operatorname{spec}(\rho)\prec\operatorname{spec}(\rho_K).
\tag{FM}\label{eq:fqe-filter-majorization}
\]
In particular, filtering a single qubit of any pure critical four-qubit
state cannot increase the R\'enyi entropy of \emph{any} of the three
balanced cuts, at any $q>0$.
\end{lemma}
\begin{proof}
The positive factor in the polar decomposition of $K$ suffices,
since the unitary factor only conjugates $\rho_K$. By a positive
rescaling assume $\operatorname{tr}K^\dagger K=2$; the given
single-qubit marginal makes the filtered state already normalised.
Then $K^\dagger K=I+s\,\mathbf n\cdot\boldsymbol\sigma$,
$0\le s\le1$, where $\mathbf n$ is a real unit vector. Put
$N=(\mathbf n\cdot\boldsymbol\sigma)\otimes I$.

Let $\Theta=(\sigma_y\otimes\sigma_y)\mathcal K$ be the
two-qubit spin-flip antiunitary, where $\mathcal K$ is
complex conjugation in the computational basis.
Expanding $\rho$ in Pauli matrices shows that
$\Theta\rho\Theta^{-1}=\rho$, because the two one-site
Bloch vectors vanish and spin flip changes the sign of
each individual Pauli factor. Also
$\Theta N\Theta^{-1}=-N$.
The matrices
\[
B_\pm=\rho^{1/2}(I\pm sN)\rho^{1/2}
\]
are positive, have the same spectrum by conjugation with
$\Theta$, and satisfy $(B_++B_-)/2=\rho$.
The nonzero eigenvalues of $B_+$ coincide with those of
$(K\otimes I)\rho(K^\dagger\otimes I)$, because these
matrices have the form $X^\dagger X$ and $XX^\dagger$.
Convexity of every Ky Fan sum under averaging therefore gives
$\operatorname{spec}(\rho)\prec\operatorname{spec}(B_+)
=\operatorname{spec}(\rho_K)$.

For a critical four-qubit state, each qubit marginal is $I/2$.
On each balanced cut choose the two-qubit side containing the
filtered site. Its reduced state has both one-qubit marginals $I/2$;
the two sides of a pure bipartition have the same nonzero spectrum.
Apply the result above and Schur concavity of R\'enyi entropy.
\end{proof}

\paragraph{Two filters need not preserve this ordering.}
The single-site conclusion is \emph{not} an orbitwise
monotonicity theorem. For an exact counterexample let
\[
z=(2+2i)(|0011\rangle+|1100\rangle)
 +|0101\rangle-|0110\rangle-|1001\rangle+|1010\rangle.
\]
It has squared norm $20$ and all four one-qubit marginals
maximally mixed. Apply $\operatorname{diag}(2,1)$ at
\emph{both} sites 2 and 3, and normalise the resulting
unnormalised vector $z'$ of squared norm $89$.
The three cut spectra of $z/\sqrt{20}$ are, respectively,
\[
(2/5,2/5,1/5,0),\qquad
(13/20,1/4,1/20,1/20),\qquad
(13/20,1/4,1/20,1/20),
\]
whereas those of $z'/\sqrt{89}$ are
\[
(32,32,25,0)/89,\quad
(65,16,4,4)/89,\quad
(52,20,16,1)/89.
\]
Consequently, their products of cut purities are
\[
\prod_jp_j(z/\sqrt{20})=\frac{21609}{250000}
\;>\;
\frac{40544579889}{496981290961}
=\prod_jp_j(z'/\sqrt{89}).
\tag{FN}\label{eq:fqe-filter-nogo}
\]
This proves that the \emph{mean} R\'enyi-2 entropy increases
under two local filters. The same is true at $q=3/2$, inside
the intermediate-order interval: if $M(\lambda)=\sum_i\lambda_i^{3/2}$,
the original moments obey $M_1>59/100$ and
$M_2=M_3>67/100$, while the filtered moments obey
$M'_1<581/1000$, $M'_2<720/1000$, and $M'_3<631/1000$.
Each strict bound follows by squaring rational bounds on the
square roots of the displayed rational eigenvalues. Thus
\[
\prod_jM_j>\frac{264851}{10^6}>
\frac{3299499}{12500000}>\prod_jM'_j,
\]
so $\mathcal E_{3/2}(z'/\sqrt{89})
>\mathcal E_{3/2}(z/\sqrt{20})$.
Both states obey the global HS--cluster envelope of Corollary~\ref{cor:fqe-renyi-gw}.
The valid single-filter lemma therefore cannot be iterated to
replace all noncritical states by critical representatives.
